\documentclass[aspectratio=169,11pt]{beamer}
\usetheme{Madrid}
\usecolortheme{seahorse}
\setbeamertemplate{navigation symbols}{}
\setbeamertemplate{headline}{}
\setbeamertemplate{footline}[frame number]
\setbeamertemplate{frametitle continuation}{}
\setbeamertemplate{frametitle}{\nointerlineskip\begin{beamercolorbox}[wd=\paperwidth,sep=0.45ex,leftskip=7.5mm,rightskip=7.5mm]{frametitle}\usebeamerfont{frametitle}\insertframetitle\end{beamercolorbox}}
\setbeamerfont{frametitle}{size=\large}
\setbeamerfont{normal text}{size=\normalsize}
\setbeamersize{text margin left=8mm,text margin right=8mm}
\setbeameroption{hide notes}
\usepackage[T1]{fontenc}
\usepackage{tabularx}
\usepackage{array}
\usepackage{hyperref}
\usepackage{fancyvrb}
\pdfstringdefDisableCommands{\def\translate#1{#1}}
\title{Mizar Evo}
\subtitle{Detailed discussion deck for the Bia\l{}ystok Mizar team}
\author{Mizar Evo project}
\date{September 2026}

\begin{document}
\begin{frame}
\titlepage
\end{frame}

\section{Part 0. Opening}
\begin{frame}[fragile,t]{0.1 - Title}
\smallskip
\textbf{Title:}
\begin{Verbatim}[fontsize=\small,frame=single,framesep=0.8mm,framerule=0.35pt,rulecolor=\color{black!45},xleftmargin=0.5mm,xrightmargin=0.5mm]
Mizar Evo
Readable, AI-Ready, Scalable Formal Mathematics
\end{Verbatim}
\smallskip
\textbf{Subtitle:}
\begin{Verbatim}[fontsize=\small,frame=single,framesep=0.8mm,framerule=0.35pt,rulecolor=\color{black!45},xleftmargin=0.5mm,xrightmargin=0.5mm]
Eight problems, eight proposals
A discussion with the Bialystok Mizar team, September 2026
\end{Verbatim}
\note{
\smallskip
\textbf{Presenter script:}
\smallskip
\textbf{Say:}
\begin{itemize}
\setlength{\itemsep}{1pt}
\setlength{\parsep}{0pt}
\setlength{\parskip}{0pt}
\item Begin with gratitude.
\item State that this is a design review with the people best positioned to judge whether the project is still recognizably Mizar.
\item Nothing here is frozen; the goal of the visit is to collect objections.
\end{itemize}
}
\end{frame}

\begin{frame}[fragile,t]{0.2 - A First Look}
\smallskip
\textbf{Legacy Mizar (exact MML excerpt):}
\begin{Verbatim}[fontsize=\footnotesize,frame=single,framesep=0.8mm,framerule=0.35pt,rulecolor=\color{black!45},xleftmargin=0.5mm,xrightmargin=0.5mm]
definition
  struct (1-sorted) addMagma (# carrier -> set, addF -> BinOp of the carrier
  #);
end;
\end{Verbatim}
\smallskip
\textbf{Mizar Evo (specification example):}
\begin{Verbatim}[fontsize=\footnotesize,frame=single,framesep=0.8mm,framerule=0.35pt,rulecolor=\color{black!45},xleftmargin=0.5mm,xrightmargin=0.5mm]
definition
  struct AddMagma where
    field carrier -> set;
    field add -> BinOp of carrier;
  end;

  inherit AddMagma extends Magma where
    field carrier from carrier;
    field add from binop;    :: renamed view
  end;
end;
\end{Verbatim}
\note{
\smallskip
\textbf{Presenter script:}
\smallskip
\textbf{Say:}
\begin{itemize}
\setlength{\itemsep}{1pt}
\setlength{\parsep}{0pt}
\setlength{\parskip}{0pt}
\item Source: current MML \texttt{algstr\_0.miz}, lines 37-40.
\item This is the whole talk in one slide: it still reads as Mizar, the mathematics is unchanged, and the things that used to be implicit - the parent link, the field mapping - are now visible, checkable source text.
\end{itemize}
}
\end{frame}

\begin{frame}[fragile,t]{0.3 - The Proposal In One Sentence}
\smallskip
\textbf{Key Phrase:}
\begin{Verbatim}[fontsize=\small,frame=single,framesep=0.8mm,framerule=0.35pt,rulecolor=\color{black!45},xleftmargin=0.5mm,xrightmargin=0.5mm]
Preserve Mizar's mathematical vernacular.
Modernize the compiler, verifier, artifact, and publication layers.
\end{Verbatim}
\smallskip
\textbf{What this talk does not claim:}
\begin{itemize}
\setlength{\itemsep}{1pt}
\setlength{\parsep}{0pt}
\setlength{\parskip}{0pt}
\item Full MML migration is not complete.
\item The final language standard is not frozen.
\item AI assistance is not a substitute for proof checking.
\item The current Mizar system's achievements are the baseline, not the problem.
\end{itemize}
\note{
\smallskip
\textbf{Presenter script:}
\smallskip
\textbf{Read aloud:}
\begin{itemize}
\setlength{\itemsep}{1pt}
\setlength{\parsep}{0pt}
\setlength{\parskip}{0pt}
\item Full MML migration is not complete.
\item The final language standard is not frozen.
\item AI assistance is not a substitute for proof checking.
\end{itemize}
After the example, say which design obligation the syntax makes visible.

}
\end{frame}

\begin{frame}[fragile,t]{0.4 - How To Read The Examples}
\smallskip
\textbf{Every code example is labeled:}
\begin{footnotesize}
\setlength{\tabcolsep}{1.5pt}
\renewcommand{\arraystretch}{1.0}
\begin{tabularx}{\textwidth}{|>{\raggedright\arraybackslash}X|>{\raggedright\arraybackslash}X|}
\hline
Label & Meaning \\
\hline
exact MML excerpt & verbatim current MML text, with article and line numbers \\
\hline
specification example & from the Mizar Evo language specification \\
\hline
sketch & illustrative; not yet fixed by the specification \\
\hline
\end{tabularx}
\end{footnotesize}

\smallskip
\textbf{Reading rule:}
\begin{itemize}
\setlength{\itemsep}{1pt}
\setlength{\parsep}{0pt}
\setlength{\parskip}{0pt}
\item No EBNF appears in this talk. The language is shown through examples only; \texttt{doc/spec/en/} remains the authority for grammar and edge cases.
\end{itemize}
\note{
\smallskip
\textbf{Presenter script:}
\smallskip
\textbf{Read aloud:}
\begin{itemize}
\setlength{\itemsep}{1pt}
\setlength{\parsep}{0pt}
\setlength{\parskip}{0pt}
\item No EBNF appears in this talk. The language is shown through examples only; doc/spec/en/ remains the authority for grammar and edge cases.
\end{itemize}
For the table, read the row labels first, then the design reason.

}
\end{frame}

\begin{frame}[fragile,t]{0.5 - What We Need From This Visit}
\smallskip
\textbf{Key Points:}
\begin{itemize}
\setlength{\itemsep}{1pt}
\setlength{\parsep}{0pt}
\setlength{\parskip}{0pt}
\item identify compatibility constraints that are invisible from outside MML work;
\item choose migration benchmark articles that are small but representative;
\item review the trust boundary for ATP search and kernel checking;
\item discuss how Formalized Mathematics should link to a package library;
\item collect the objections we have not thought of.
\end{itemize}
\smallskip
\textbf{Running question:}
\begin{Verbatim}[fontsize=\small,frame=single,framesep=0.8mm,framerule=0.35pt,rulecolor=\color{black!45},xleftmargin=0.5mm,xrightmargin=0.5mm]
What must Mizar Evo preserve so that the Mizar community
still recognizes it as Mizar?
\end{Verbatim}
\note{
\smallskip
\textbf{Presenter script:}
\smallskip
\textbf{Read aloud:}
\begin{itemize}
\setlength{\itemsep}{1pt}
\setlength{\parsep}{0pt}
\setlength{\parskip}{0pt}
\item identify compatibility constraints that are invisible from outside MML work;
\item choose migration benchmark articles that are small but representative;
\item review the trust boundary for ATP search and kernel checking;
\end{itemize}
After the example, say which design obligation the syntax makes visible.

}
\end{frame}

\section{Part 1. Why Now}
\begin{frame}[fragile,t]{1.1 - What Mizar Got Right}
\smallskip
\textbf{Key Points:}
\begin{itemize}
\setlength{\itemsep}{1pt}
\setlength{\parsep}{0pt}
\setlength{\parskip}{0pt}
\item declarative proof text that reads as mathematics;
\item soft types, modes, and adjective-rich vocabulary;
\item attributes, registrations, and clusters as reusable automation;
\item a mature curated library (MML 5.94.1493: 1493 articles);
\item a publication culture around formal articles (Formalized Mathematics).
\end{itemize}
\smallskip
\textbf{Message:}
\begin{itemize}
\setlength{\itemsep}{1pt}
\setlength{\parsep}{0pt}
\setlength{\parskip}{0pt}
\item These strengths are the design baseline. Every proposal in this talk is judged by whether it protects them.
\end{itemize}
\note{
\smallskip
\textbf{Presenter script:}
\smallskip
\textbf{Say:}
\begin{itemize}
\setlength{\itemsep}{1pt}
\setlength{\parsep}{0pt}
\setlength{\parskip}{0pt}
\item Do not lecture the audience about their own system; this frame is one minute of shared ground, not a tutorial.
\end{itemize}
\smallskip
\textbf{Read aloud:}
\begin{itemize}
\setlength{\itemsep}{1pt}
\setlength{\parsep}{0pt}
\setlength{\parskip}{0pt}
\item declarative proof text that reads as mathematics;
\item soft types, modes, and adjective-rich vocabulary;
\item attributes, registrations, and clusters as reusable automation;
\end{itemize}
}
\end{frame}

\begin{frame}[fragile,t]{1.2 - Pressure One: Scale}
\smallskip
\textbf{Key Points:}
\begin{itemize}
\setlength{\itemsep}{1pt}
\setlength{\parsep}{0pt}
\setlength{\parskip}{0pt}
\item MML has grown to roughly 1500 interdependent articles.
\item The unit of dependency, review, and reuse is the whole article.
\item Article environments are resolved by tooling, but the resolved dependency surface is not visible in the source a human reviews.
\item Whole-library maintenance operations (renames, refactorings, revisions) carry risk that grows with library size.
\end{itemize}
\smallskip
\textbf{Message:}
\begin{itemize}
\setlength{\itemsep}{1pt}
\setlength{\parsep}{0pt}
\setlength{\parskip}{0pt}
\item The problem is not that current Mizar is wrong. It is that article-level boundaries were designed for a smaller library.
\end{itemize}
\note{
\smallskip
\textbf{Presenter script:}
\smallskip
\textbf{Read aloud:}
\begin{itemize}
\setlength{\itemsep}{1pt}
\setlength{\parsep}{0pt}
\setlength{\parskip}{0pt}
\item MML has grown to roughly 1500 interdependent articles.
\item The unit of dependency, review, and reuse is the whole article.
\item Article environments are resolved by tooling, but the resolved dependency surface is not visible in the source a human reviews.
\item Whole-library maintenance operations (renames, refactorings, revisions) carry risk that grows with library size.
\item The problem is not that current Mizar is wrong. It is that article-level boundaries were designed for a smaller library.
\end{itemize}
}
\end{frame}

\begin{frame}[fragile,t]{1.3 - Pressure Two: Tooling Expectations}
\smallskip
\textbf{Key Points:}
\begin{itemize}
\setlength{\itemsep}{1pt}
\setlength{\parsep}{0pt}
\setlength{\parskip}{0pt}
\item Editors are expected to give instant, partial, resilient feedback.
\item Builds are expected to be reproducible from a manifest and lockfile.
\item Reuse is expected to work at package granularity with versioning.
\item Documentation is expected to be generated, linked, and browsable.
\end{itemize}
\smallskip
\textbf{Message:}
\begin{itemize}
\setlength{\itemsep}{1pt}
\setlength{\parsep}{0pt}
\setlength{\parskip}{0pt}
\item These expectations were set by mainstream language ecosystems; new users arrive with them, and formal libraries are judged against them.
\end{itemize}
\note{
\smallskip
\textbf{Presenter script:}
\smallskip
\textbf{Read aloud:}
\begin{itemize}
\setlength{\itemsep}{1pt}
\setlength{\parsep}{0pt}
\setlength{\parskip}{0pt}
\item Editors are expected to give instant, partial, resilient feedback.
\item Builds are expected to be reproducible from a manifest and lockfile.
\item Reuse is expected to work at package granularity with versioning.
\item Documentation is expected to be generated, linked, and browsable.
\item These expectations were set by mainstream language ecosystems; new users arrive with them, and formal libraries are judged against them.
\end{itemize}
}
\end{frame}

\begin{frame}[fragile,t]{1.4 - Pressure Three: AI}
\smallskip
\textbf{Key Points:}
\begin{itemize}
\setlength{\itemsep}{1pt}
\setlength{\parsep}{0pt}
\setlength{\parskip}{0pt}
\item AI agents are already useful for search, explanation, and repair.
\item They need bounded, structured, source-anchored context - not a dump of the whole library.
\item Their output must never define proof truth; verification must stay independent of the strength of the assistant.
\item Readable source is an advantage here: stable local text patterns are what AI edits and retrieval work best on.
\end{itemize}
\smallskip
\textbf{Message:}
\begin{itemize}
\setlength{\itemsep}{1pt}
\setlength{\parsep}{0pt}
\setlength{\parskip}{0pt}
\item Mizar's readability is not a nostalgic asset. It is exactly what makes safe AI assistance possible.
\end{itemize}
\note{
\smallskip
\textbf{Presenter script:}
\smallskip
\textbf{Read aloud:}
\begin{itemize}
\setlength{\itemsep}{1pt}
\setlength{\parsep}{0pt}
\setlength{\parskip}{0pt}
\item AI agents are already useful for search, explanation, and repair.
\item They need bounded, structured, source-anchored context - not a dump of the whole library.
\item Their output must never define proof truth; verification must stay independent of the strength of the assistant.
\item Readable source is an advantage here: stable local text patterns are what AI edits and retrieval work best on.
\item Mizar's readability is not a nostalgic asset. It is exactly what makes safe AI assistance possible.
\end{itemize}
}
\end{frame}

\begin{frame}[fragile,t]{1.5 - The Design Rule}
\smallskip
\textbf{Key Phrase:}
\begin{Verbatim}[fontsize=\small,frame=single,framesep=0.8mm,framerule=0.35pt,rulecolor=\color{black!45},xleftmargin=0.5mm,xrightmargin=0.5mm]
Do not trade away readability to gain automation.
Use automation to protect and extend readability.
\end{Verbatim}
\smallskip
\textbf{Three pillars, one test:}
\begin{footnotesize}
\setlength{\tabcolsep}{1.5pt}
\renewcommand{\arraystretch}{1.0}
\begin{tabularx}{\textwidth}{|>{\raggedright\arraybackslash}X|>{\raggedright\arraybackslash}X|}
\hline
Pillar & Test for every feature \\
\hline
Readability & does proof text still read as mathematics? \\
\hline
AI-readiness & can a tool see bounded, auditable context? \\
\hline
Scalability & do boundaries stay stable as the library grows? \\
\hline
\end{tabularx}
\end{footnotesize}

\note{
\smallskip
\textbf{Presenter script:}
\smallskip
\textbf{Say:}
\begin{itemize}
\setlength{\itemsep}{1pt}
\setlength{\parsep}{0pt}
\setlength{\parskip}{0pt}
\item The eight stories that follow each apply this rule to one concrete pain.
\end{itemize}
For the table, read the row labels first, then the design reason.

}
\end{frame}

\section{Part 2. Story 1: Dependencies You Can See}
\begin{frame}[fragile,t]{2.1 - The Pain}
\smallskip
\textbf{Legacy Mizar (exact MML excerpt):}
\begin{Verbatim}[fontsize=\footnotesize,frame=single,framesep=0.8mm,framerule=0.35pt,rulecolor=\color{black!45},xleftmargin=0.5mm,xrightmargin=0.5mm]
environ

 vocabularies XBOOLE_0, SUBSET_1, BINOP_1, ZFMISC_1, STRUCT_0, ARYTM_3,
      FUNCT_1, FUNCT_5, SUPINF_2, ARYTM_1, RELAT_1, MESFUNC1, ALGSTR_0, CARD_1;
 notations TARSKI, XBOOLE_0, SUBSET_1, ZFMISC_1, BINOP_1, FUNCT_5, ORDINAL1,
      CARD_1, STRUCT_0;
 constructors BINOP_1, STRUCT_0, ZFMISC_1, FUNCT_5;
 registrations ZFMISC_1, CARD_1, STRUCT_0;
 theorems STRUCT_0;

begin :: Additive structures
\end{Verbatim}
\note{
\smallskip
\textbf{Presenter script:}
\smallskip
\textbf{Say:}
\begin{itemize}
\setlength{\itemsep}{1pt}
\setlength{\parsep}{0pt}
\setlength{\parskip}{0pt}
\item Source: current MML \texttt{algstr\_0.miz}, lines 15-25.
\item Everyone in the room has edited one of these blocks by trial and error.
\item The point is not that \texttt{environ} is bad; it successfully drove decades of library growth. The point is what it costs at today's scale.
\end{itemize}
}
\end{frame}

\begin{frame}[fragile,t]{2.2 - Why It Hurts}
\smallskip
\textbf{Key Points:}
\begin{itemize}
\setlength{\itemsep}{1pt}
\setlength{\parsep}{0pt}
\setlength{\parskip}{0pt}
\item One symbol's origin is spread over several role lists; a reviewer cannot see which article contributes which notation, constructor, or cluster.
\item The Accommodator resolves the environment, but the resolved surface is not reviewable source text.
\item Tools cannot cache or invalidate at a finer granularity than the article.
\item Moving a theorem between articles risks breaking unknown dependents.
\end{itemize}
\smallskip
\textbf{Message:}
\begin{itemize}
\setlength{\itemsep}{1pt}
\setlength{\parsep}{0pt}
\setlength{\parskip}{0pt}
\item Implicit dependency surfaces are a fixed cost on every edit, every review, and every tool, and the cost grows with the library.
\end{itemize}
\note{
\smallskip
\textbf{Presenter script:}
\smallskip
\textbf{Read aloud:}
\begin{itemize}
\setlength{\itemsep}{1pt}
\setlength{\parsep}{0pt}
\setlength{\parskip}{0pt}
\item One symbol's origin is spread over several role lists; a reviewer cannot see which article contributes which notation, constructor, or cluster.
\item The Accommodator resolves the environment, but the resolved surface is not reviewable source text.
\item Tools cannot cache or invalidate at a finer granularity than the article.
\item Moving a theorem between articles risks breaking unknown dependents.
\item Implicit dependency surfaces are a fixed cost on every edit, every review, and every tool, and the cost grows with the library.
\end{itemize}
}
\end{frame}

\begin{frame}[fragile,t]{2.3 - The Evo Answer: Import Prelude}
\smallskip
\textbf{Mizar Evo (specification example):}
\begin{Verbatim}[fontsize=\footnotesize,frame=single,framesep=0.8mm,framerule=0.35pt,rulecolor=\color{black!45},xleftmargin=0.5mm,xrightmargin=0.5mm]
import .function;
import mml.algebra.structure.sorted;

definition
  let S be 1-sorted;
  mode BinOpDef: BinOp of S is
    Function of [: S.carrier, S.carrier :], S.carrier;
end;
\end{Verbatim}
\smallskip
\textbf{Rules that make this deterministic:}
\begin{itemize}
\setlength{\itemsep}{1pt}
\setlength{\parsep}{0pt}
\setlength{\parskip}{0pt}
\item all imports appear before the first item; no mid-file environment changes;
\item imports seed the active lexicon, and every symbol, notation, registration, and theorem is traceable to exactly one import;
\item stable fully-qualified names are derived from package and module paths.
\end{itemize}
\note{
\smallskip
\textbf{Presenter script:}
\smallskip
\textbf{Read aloud:}
\begin{itemize}
\setlength{\itemsep}{1pt}
\setlength{\parsep}{0pt}
\setlength{\parskip}{0pt}
\item all imports appear before the first item; no mid-file environment changes;
\item imports seed the active lexicon, and every symbol, notation, registration, and theorem is traceable to exactly one import;
\item stable fully-qualified names are derived from package and module paths.
\end{itemize}
After the example, say which design obligation the syntax makes visible.

}
\end{frame}

\begin{frame}[fragile,t]{2.4 - The Evo Answer: Packages}
\smallskip
\textbf{Mizar Evo (specification example):}
\begin{Verbatim}[fontsize=\footnotesize,frame=single,framesep=0.8mm,framerule=0.35pt,rulecolor=\color{black!45},xleftmargin=0.5mm,xrightmargin=0.5mm]
[package]
name    = "algebra"
version = "2.3.1"
edition = "2025"

[dependencies]
mml_core = "^1.0"
topology = { version = "^0.9", features = ["metric"] }
\end{Verbatim}
\smallskip
\textbf{Key Points:}
\begin{itemize}
\setlength{\itemsep}{1pt}
\setlength{\parsep}{0pt}
\setlength{\parskip}{0pt}
\item a manifest plus a lockfile makes every build reproducible;
\item versioned reuse (SemVer) replaces ad hoc copying between article sets.
\end{itemize}
\note{
\smallskip
\textbf{Presenter script:}
\smallskip
\textbf{Read aloud:}
\begin{itemize}
\setlength{\itemsep}{1pt}
\setlength{\parsep}{0pt}
\setlength{\parskip}{0pt}
\item a manifest plus a lockfile makes every build reproducible;
\item versioned reuse (SemVer) replaces ad hoc copying between article sets.
\end{itemize}
After the example, say which design obligation the syntax makes visible.

}
\end{frame}

\begin{frame}[fragile,t]{2.5 - Migrating The Environment}
\smallskip
\textbf{Migration map:}
\begin{footnotesize}
\setlength{\tabcolsep}{1.5pt}
\renewcommand{\arraystretch}{1.0}
\begin{tabularx}{\textwidth}{|>{\raggedright\arraybackslash}X|>{\raggedright\arraybackslash}X|}
\hline
Current environ role & Evo target \\
\hline
vocabularies & exported symbols and lexical metadata \\
\hline
notations & imported notation metadata \\
\hline
constructors & visible definitions and constructors \\
\hline
registrations & import-scoped registration index \\
\hline
requirements & package or prelude policy \\
\hline
article theorem labels & module-qualified theorem identities \\
\hline
\end{tabularx}
\end{footnotesize}

\smallskip
\textbf{Message:}
\begin{itemize}
\setlength{\itemsep}{1pt}
\setlength{\parsep}{0pt}
\setlength{\parskip}{0pt}
\item this is not a mechanical rename: current environments mix semantic, syntactic, and automation-facing roles, and migration reports must explain what each imported module actually contributes.
\end{itemize}
\note{
\smallskip
\textbf{Presenter script:}
\smallskip
\textbf{Read aloud:}
\begin{itemize}
\setlength{\itemsep}{1pt}
\setlength{\parsep}{0pt}
\setlength{\parskip}{0pt}
\item this is not a mechanical rename: current environments mix semantic, syntactic, and automation-facing roles, and migration reports must explain what each imported module actually contributes.
\end{itemize}
For the table, read the row labels first, then the design reason.

}
\end{frame}

\begin{frame}[fragile,t]{2.6 - What Is Preserved, What We Ask}
\smallskip
\textbf{Preserved:}
\begin{itemize}
\setlength{\itemsep}{1pt}
\setlength{\parsep}{0pt}
\setlength{\parskip}{0pt}
\item the mathematics and theorem identities are untouched;
\item article-style authorship remains; a module is still a readable text;
\item migration keeps origin metadata (story 8 returns to this).
\end{itemize}
\smallskip
\textbf{Questions for Bialystok:}
\begin{itemize}
\setlength{\itemsep}{1pt}
\setlength{\parsep}{0pt}
\setlength{\parskip}{0pt}
\item Which \texttt{environ} roles must remain visually familiar during migration?
\item Which current article dependencies are hardest to explain, and would make the best test cases for generated dependency reports?
\end{itemize}
\note{
\smallskip
\textbf{Presenter script:}
\smallskip
\textbf{Read aloud:}
\begin{itemize}
\setlength{\itemsep}{1pt}
\setlength{\parsep}{0pt}
\setlength{\parskip}{0pt}
\item the mathematics and theorem identities are untouched;
\item article-style authorship remains; a module is still a readable text;
\item migration keeps origin metadata (story 8 returns to this).
\item Which environ roles must remain visually familiar during migration?
\item Which current article dependencies are hardest to explain, and would make the best test cases for generated dependency reports?
\end{itemize}
}
\end{frame}

\section{Part 3. Story 2: Structures Without Hidden Merges}
\begin{frame}[fragile,t]{3.1 - The Pain}
\smallskip
\textbf{Legacy Mizar (exact MML excerpt):}
\begin{Verbatim}[fontsize=\footnotesize,frame=single,framesep=0.8mm,framerule=0.35pt,rulecolor=\color{black!45},xleftmargin=0.5mm,xrightmargin=0.5mm]
definition
  struct (1-sorted) addMagma (# carrier -> set, addF -> BinOp of the carrier
  #);
end;
\end{Verbatim}
\smallskip
\textbf{Key Points:}
\begin{itemize}
\setlength{\itemsep}{1pt}
\setlength{\parsep}{0pt}
\setlength{\parskip}{0pt}
\item parent link, fields, and selector layout are one compact declaration;
\item with multiple parents, the merge of inherited fields is implicit;
\item renamed views (additive vs multiplicative) rest on naming conventions;
\item stored data and canonical values (a zero, a unit) are not distinguished.
\end{itemize}
\note{
\smallskip
\textbf{Presenter script:}
\smallskip
\textbf{Say:}
\begin{itemize}
\setlength{\itemsep}{1pt}
\setlength{\parsep}{0pt}
\setlength{\parskip}{0pt}
\item Source: current MML \texttt{algstr\_0.miz}, lines 37-40.
\item Acknowledge that this compactness was a feature: structures stayed close to informal mathematical writing.
\end{itemize}
\smallskip
\textbf{Read aloud:}
\begin{itemize}
\setlength{\itemsep}{1pt}
\setlength{\parsep}{0pt}
\setlength{\parskip}{0pt}
\item parent link, fields, and selector layout are one compact declaration;
\item with multiple parents, the merge of inherited fields is implicit;
\item renamed views (additive vs multiplicative) rest on naming conventions;
\end{itemize}
}
\end{frame}

\begin{frame}[fragile,t]{3.2 - Why It Hurts}
\smallskip
\textbf{Key Points:}
\begin{itemize}
\setlength{\itemsep}{1pt}
\setlength{\parsep}{0pt}
\setlength{\parskip}{0pt}
\item Diamond inheritance (one structure reachable through two parent paths) is resolved by convention and declaration order, not by checkable source.
\item A migration tool cannot ask "is this selector intrinsic data, a canonical value with obligations, or an inherited view?" - the syntax does not say.
\item Errors surface far from their cause, as type mismatches in later articles.
\end{itemize}
\smallskip
\textbf{Message:}
\begin{itemize}
\setlength{\itemsep}{1pt}
\setlength{\parsep}{0pt}
\setlength{\parskip}{0pt}
\item At MML scale, structure inheritance is a graph maintenance problem, and the graph deserves explicit, checkable edges.
\end{itemize}
\note{
\smallskip
\textbf{Presenter script:}
\smallskip
\textbf{Read aloud:}
\begin{itemize}
\setlength{\itemsep}{1pt}
\setlength{\parsep}{0pt}
\setlength{\parskip}{0pt}
\item Diamond inheritance (one structure reachable through two parent paths) is resolved by convention and declaration order, not by checkable source.
\item A migration tool cannot ask "is this selector intrinsic data, a canonical value with obligations, or an inherited view?" - the syntax does not say.
\item Errors surface far from their cause, as type mismatches in later articles.
\item At MML scale, structure inheritance is a graph maintenance problem, and the graph deserves explicit, checkable edges.
\end{itemize}
}
\end{frame}

\begin{frame}[fragile,t]{3.3 - The Evo Answer: Field, Property, Attribute}
\smallskip
\textbf{Mizar Evo (specification example):}
\begin{Verbatim}[fontsize=\small,frame=single,framesep=0.8mm,framerule=0.35pt,rulecolor=\color{black!45},xleftmargin=0.5mm,xrightmargin=0.5mm]
definition
  struct AddLoopStr where
    field carrier -> set;
    field add -> BinOp of carrier;
    property zero -> Element of carrier;
  end;
end;
\end{Verbatim}
\smallskip
\textbf{Three distinct concepts:}
\begin{footnotesize}
\setlength{\tabcolsep}{1.5pt}
\renewcommand{\arraystretch}{1.0}
\begin{tabularx}{\textwidth}{|>{\raggedright\arraybackslash}X|>{\raggedright\arraybackslash}X|>{\raggedright\arraybackslash}X|}
\hline
Concept & Meaning & Consequence \\
\hline
\texttt{field} & intrinsic data supplied by the value & constructor shape, equality \\
\hline
\texttt{property} & uniquely determined canonical value & existence/uniqueness obligations \\
\hline
\texttt{attribute} & predicate-style refinement & cluster propagation, not layout \\
\hline
\end{tabularx}
\end{footnotesize}

\note{
\smallskip
\textbf{Presenter script:}
After the example, say which design obligation the syntax makes visible.

For the table, read the row labels first, then the design reason.

Use this as the transition: The Evo Answer: Field, Property, Attribute. Point to the visible example, question, or diagram before moving on.

}
\end{frame}

\begin{frame}[fragile,t]{3.4 - The Evo Answer: Explicit Inheritance}
\smallskip
\textbf{Mizar Evo (specification example):}
\begin{Verbatim}[fontsize=\footnotesize,frame=single,framesep=0.8mm,framerule=0.35pt,rulecolor=\color{black!45},xleftmargin=0.5mm,xrightmargin=0.5mm]
definition
  struct AddMagma where
    field carrier -> set;
    field add -> BinOp of carrier;
  end;

  inherit AddMagma extends Magma where
    field carrier from carrier;
    field add from binop;    :: renamed
  end;
end;
\end{Verbatim}
\smallskip
\textbf{Key Points:}
\begin{itemize}
\setlength{\itemsep}{1pt}
\setlength{\parsep}{0pt}
\setlength{\parskip}{0pt}
\item one \texttt{inherit} statement per parent; mapping and renaming are source text;
\item narrowing and type conversions carry \texttt{coherence} proof obligations;
\item the additive/multiplicative naming convention becomes a checked view.
\end{itemize}
\note{
\smallskip
\textbf{Presenter script:}
\smallskip
\textbf{Read aloud:}
\begin{itemize}
\setlength{\itemsep}{1pt}
\setlength{\parsep}{0pt}
\setlength{\parskip}{0pt}
\item one inherit statement per parent; mapping and renaming are source text;
\item narrowing and type conversions carry coherence proof obligations;
\item the additive/multiplicative naming convention becomes a checked view.
\end{itemize}
After the example, say which design obligation the syntax makes visible.

}
\end{frame}

\begin{frame}[fragile,t]{3.5 - The Evo Answer: Diamonds Become Checkable}
\smallskip
\textbf{Mizar Evo (specification example):}
\begin{Verbatim}[fontsize=\footnotesize,frame=single,framesep=0.8mm,framerule=0.35pt,rulecolor=\color{black!45},xleftmargin=0.5mm,xrightmargin=0.5mm]
struct DoubleLoopStr where
  field carrier -> set;
  field add -> BinOp of carrier;
  field mul -> BinOp of carrier;
  property zero -> Element of carrier;
  property one -> Element of carrier;
end;

inherit DoubleLoopStr extends AddLoopStr;
inherit DoubleLoopStr extends MulLoopStr;
\end{Verbatim}
\smallskip
\textbf{Message:}
\begin{itemize}
\setlength{\itemsep}{1pt}
\setlength{\parsep}{0pt}
\setlength{\parskip}{0pt}
\item The analyzer must check that both inheritance paths introduce the same components: \texttt{add -> LoopStr.binop -> Magma.binop} along path one must agree with \texttt{add -> AddMagma.add -> Magma.binop} along path two.
\item Diamond inheritance becomes a diagnostic with source spans, not a silent merge decided by declaration order.
\end{itemize}
\note{
\smallskip
\textbf{Presenter script:}
\smallskip
\textbf{Read aloud:}
\begin{itemize}
\setlength{\itemsep}{1pt}
\setlength{\parsep}{0pt}
\setlength{\parskip}{0pt}
\item The analyzer must check that both inheritance paths introduce the same components: add -> LoopStr.binop -> Magma.binop along path one must agree with add -> AddMagma.add -> Magma.binop along path two.
\item Diamond inheritance becomes a diagnostic with source spans, not a silent merge decided by declaration order.
\end{itemize}
After the example, say which design obligation the syntax makes visible.

}
\end{frame}

\begin{frame}[fragile,t]{3.6 - What Is Preserved, What We Ask}
\smallskip
\textbf{Preserved:}
\begin{itemize}
\setlength{\itemsep}{1pt}
\setlength{\parsep}{0pt}
\setlength{\parskip}{0pt}
\item structures remain Mizar structures: carriers, selectors, \texttt{Element of};
\item aggregate compactness is traded only where a hidden decision existed.
\end{itemize}
\smallskip
\textbf{Questions for Bialystok:}
\begin{itemize}
\setlength{\itemsep}{1pt}
\setlength{\parsep}{0pt}
\setlength{\parskip}{0pt}
\item Is the \texttt{field} / \texttt{property} / \texttt{attribute} split readable in real algebraic articles, or does it over-annotate simple cases?
\item Is one-parent-per-\texttt{inherit} acceptable for the inheritance-heavy parts of MML, such as the \texttt{ALGSTR} and topology hierarchies?
\item Which MML structures would be the best diamond test cases?
\end{itemize}
\note{
\smallskip
\textbf{Presenter script:}
\smallskip
\textbf{Read aloud:}
\begin{itemize}
\setlength{\itemsep}{1pt}
\setlength{\parsep}{0pt}
\setlength{\parskip}{0pt}
\item structures remain Mizar structures: carriers, selectors, Element of;
\item aggregate compactness is traded only where a hidden decision existed.
\item Is the field / property / attribute split readable in real algebraic articles, or does it over-annotate simple cases?
\item Is one-parent-per-inherit acceptable for the inheritance-heavy parts of MML, such as the ALGSTR and topology hierarchies?
\item Which MML structures would be the best diamond test cases?
\end{itemize}
}
\end{frame}

\section{Part 4. Story 3: Automation You Can Audit}
\begin{frame}[fragile,t]{4.1 - The Pain}
\smallskip
\textbf{Legacy Mizar (exact MML excerpt):}
\begin{Verbatim}[fontsize=\small,frame=single,framesep=0.8mm,framerule=0.35pt,rulecolor=\color{black!45},xleftmargin=0.5mm,xrightmargin=0.5mm]
registration
  let M be addMagma;
  cluster right_add-cancelable left_add-cancelable -> add-cancelable for
Element
    of M;
  coherence;
end;
\end{Verbatim}
\note{
\smallskip
\textbf{Presenter script:}
\smallskip
\textbf{Say:}
\begin{itemize}
\setlength{\itemsep}{1pt}
\setlength{\parsep}{0pt}
\setlength{\parskip}{0pt}
\item Source: current MML \texttt{algstr\_0.miz}, lines 104-109.
\item Registrations are one of Mizar's best ideas: adjectives propagate silently and proofs stay short. The pain is not the mechanism but its opacity.
\end{itemize}
}
\end{frame}

\begin{frame}[fragile,t]{4.2 - Why It Hurts}
\smallskip
\textbf{Key Points:}
\begin{itemize}
\setlength{\itemsep}{1pt}
\setlength{\parsep}{0pt}
\setlength{\parskip}{0pt}
\item When a proof fails, "why does the checker not see that this is a Group?" has no local answer; the cause lives somewhere in the environment.
\item Which registrations fired, in which order, is invisible; the automation is powerful, but its explanation does not scale with its power.
\item For an AI assistant the situation is worse: it must guess the cluster state instead of reading it.
\end{itemize}
\smallskip
\textbf{Message:}
\begin{itemize}
\setlength{\itemsep}{1pt}
\setlength{\parsep}{0pt}
\setlength{\parskip}{0pt}
\item Automation that cannot explain itself becomes a maintenance liability at library scale - even when it is sound.
\end{itemize}
\note{
\smallskip
\textbf{Presenter script:}
\smallskip
\textbf{Read aloud:}
\begin{itemize}
\setlength{\itemsep}{1pt}
\setlength{\parsep}{0pt}
\setlength{\parskip}{0pt}
\item When a proof fails, "why does the checker not see that this is a Group?" has no local answer; the cause lives somewhere in the environment.
\item Which registrations fired, in which order, is invisible; the automation is powerful, but its explanation does not scale with its power.
\item For an AI assistant the situation is worse: it must guess the cluster state instead of reading it.
\item Automation that cannot explain itself becomes a maintenance liability at library scale - even when it is sound.
\end{itemize}
}
\end{frame}

\begin{frame}[fragile,t]{4.3 - The Evo Answer: Labeled, Traceable Registrations}
\smallskip
\textbf{Mizar Evo (specification example):}
\begin{Verbatim}[fontsize=\small,frame=single,framesep=0.8mm,framerule=0.35pt,rulecolor=\color{black!45},xleftmargin=0.5mm,xrightmargin=0.5mm]
registration
  cluster EmptyImpliesFinite: empty -> finite for set;
  coherence proof ... end;

  cluster FiniteImpliesCountable: finite -> countable for set;
  coherence proof ... end;
end;
\end{Verbatim}
\smallskip
\textbf{Key Points:}
\begin{itemize}
\setlength{\itemsep}{1pt}
\setlength{\parsep}{0pt}
\setlength{\parskip}{0pt}
\item every registration item carries a required label: citable in \texttt{by}, reported in diagnostics, part of the module interface;
\item the verifier stores an import-filtered cluster resolution graph;
\item \texttt{@show\_resolution} and explanation artifacts answer "why (not)?" with the actual chain, e.g. empty -> finite -> countable.
\end{itemize}
\note{
\smallskip
\textbf{Presenter script:}
\smallskip
\textbf{Read aloud:}
\begin{itemize}
\setlength{\itemsep}{1pt}
\setlength{\parsep}{0pt}
\setlength{\parskip}{0pt}
\item every registration item carries a required label: citable in by, reported in diagnostics, part of the module interface;
\item the verifier stores an import-filtered cluster resolution graph;
\item @show\_resolution and explanation artifacts answer "why (not)?" with the actual chain, e.g. empty -> finite -> countable.
\end{itemize}
After the example, say which design obligation the syntax makes visible.

}
\end{frame}

\begin{frame}[fragile,t]{4.4 - The Evo Answer: Oriented Reductions}
\smallskip
\textbf{Mizar Evo (specification example):}
\begin{Verbatim}[fontsize=\footnotesize,frame=single,framesep=0.8mm,framerule=0.35pt,rulecolor=\color{black!45},xleftmargin=0.5mm,xrightmargin=0.5mm]
registration
  let n be Nat;
  reduce NatAddZero: n + 0 to n;
  reducibility
  proof
    let n be Nat;
    thus n + 0 = n by mml.number.natural.Nat_add_zero;
  end;
end;
\end{Verbatim}
\smallskip
\textbf{Key Points:}
\begin{itemize}
\setlength{\itemsep}{1pt}
\setlength{\parsep}{0pt}
\setlength{\parskip}{0pt}
\item a reduction is an oriented simplification backed by an equality proof;
\item the right side must be strictly smaller, so imported rules cannot loop, and rule selection is deterministic (specificity first, FQN tie-break);
\item unoriented identification idioms become auditable \texttt{reduce} items.
\end{itemize}
\note{
\smallskip
\textbf{Presenter script:}
\smallskip
\textbf{Read aloud:}
\begin{itemize}
\setlength{\itemsep}{1pt}
\setlength{\parsep}{0pt}
\setlength{\parskip}{0pt}
\item a reduction is an oriented simplification backed by an equality proof;
\item the right side must be strictly smaller, so imported rules cannot loop, and rule selection is deterministic (specificity first, FQN tie-break);
\item unoriented identification idioms become auditable reduce items.
\end{itemize}
After the example, say which design obligation the syntax makes visible.

}
\end{frame}

\begin{frame}[fragile,t]{4.5 - What Is Preserved, What We Ask}
\smallskip
\textbf{Preserved:}
\begin{itemize}
\setlength{\itemsep}{1pt}
\setlength{\parsep}{0pt}
\setlength{\parskip}{0pt}
\item registrations and clusters remain first-class; proofs stay short;
\item no new proof text is required at use sites - only traces are added.
\end{itemize}
\smallskip
\textbf{Questions for Bialystok:}
\begin{itemize}
\setlength{\itemsep}{1pt}
\setlength{\parsep}{0pt}
\setlength{\parskip}{0pt}
\item Which cluster explanations would most reduce daily friction: failure explanations, firing traces, or difference reports between environments?
\item Which MML article families stress registrations hardest and should become migration benchmarks for the cluster graph?
\end{itemize}
\note{
\smallskip
\textbf{Presenter script:}
\smallskip
\textbf{Read aloud:}
\begin{itemize}
\setlength{\itemsep}{1pt}
\setlength{\parsep}{0pt}
\setlength{\parskip}{0pt}
\item registrations and clusters remain first-class; proofs stay short;
\item no new proof text is required at use sites - only traces are added.
\item Which cluster explanations would most reduce daily friction: failure explanations, firing traces, or difference reports between environments?
\item Which MML article families stress registrations hardest and should become migration benchmarks for the cluster graph?
\end{itemize}
}
\end{frame}

\section{Part 5. Story 4: Powerful Search, Small Trust}
\begin{frame}[fragile,t]{5.1 - The Pain}
\smallskip
\textbf{Key Points:}
\begin{itemize}
\setlength{\itemsep}{1pt}
\setlength{\parsep}{0pt}
\setlength{\parskip}{0pt}
\item Users want stronger automation: bigger \texttt{by} steps, hammer-style search.
\item But in a monolithic verifier, every gain in search power grows the code that must be trusted.
\item External provers (ATPs) are strong exactly where Mizar's core is first-order - and they are the least auditable component of all.
\end{itemize}
\smallskip
\textbf{Key Phrase:}
\begin{Verbatim}[fontsize=\small,frame=single,framesep=0.8mm,framerule=0.35pt,rulecolor=\color{black!45},xleftmargin=0.5mm,xrightmargin=0.5mm]
How do we get modern proof search
without trusting the searcher?
\end{Verbatim}
\note{
\smallskip
\textbf{Presenter script:}
\smallskip
\textbf{Read aloud:}
\begin{itemize}
\setlength{\itemsep}{1pt}
\setlength{\parsep}{0pt}
\setlength{\parskip}{0pt}
\item Users want stronger automation: bigger by steps, hammer-style search.
\item But in a monolithic verifier, every gain in search power grows the code that must be trusted.
\item External provers (ATPs) are strong exactly where Mizar's core is first-order - and they are the least auditable component of all.
\end{itemize}
After the example, say which design obligation the syntax makes visible.

}
\end{frame}

\begin{frame}[fragile,t]{5.2 - The Evo Answer: A Reasoning Boundary}
\begin{Verbatim}[fontsize=\small,frame=single,framesep=0.8mm,framerule=0.35pt,rulecolor=\color{black!45},xleftmargin=0.5mm,xrightmargin=0.5mm]
Mizar-side semantics
  -> well-typed, resolved obligations
ATP-side search
  -> candidate proof evidence
kernel-side checking
  -> accepted or rejected proof status
\end{Verbatim}
\smallskip
\textbf{Key rules:}
\begin{itemize}
\setlength{\itemsep}{1pt}
\setlength{\parsep}{0pt}
\setlength{\parskip}{0pt}
\item ATPs never resolve names, infer types, expand clusters, or pick overloads;
\item the kernel never searches for proofs;
\item deterministic pre-ATP discharge needs replayable evidence too - nothing is accepted because an earlier phase said "done".
\end{itemize}
\note{
\smallskip
\textbf{Presenter script:}
\smallskip
\textbf{Read aloud:}
\begin{itemize}
\setlength{\itemsep}{1pt}
\setlength{\parsep}{0pt}
\setlength{\parskip}{0pt}
\item ATPs never resolve names, infer types, expand clusters, or pick overloads;
\item the kernel never searches for proofs;
\item deterministic pre-ATP discharge needs replayable evidence too - nothing is accepted because an earlier phase said "done".
\end{itemize}
After the example, say which design obligation the syntax makes visible.

}
\end{frame}

\begin{frame}[fragile,t]{5.3 - Certificates, Not Success Bits}
\begin{Verbatim}[fontsize=\footnotesize,frame=single,framesep=0.8mm,framerule=0.35pt,rulecolor=\color{black!45},xleftmargin=0.5mm,xrightmargin=0.5mm]
Certificate
  target VC fingerprint
  kernel profile
  imported facts and hashes
  generated clauses
  substitutions
  resolution trace
  final goal
\end{Verbatim}
\smallskip
\textbf{Key Points:}
\begin{itemize}
\setlength{\itemsep}{1pt}
\setlength{\parsep}{0pt}
\setlength{\parskip}{0pt}
\item accepted evidence is normalized into replayable certificate data;
\item the kernel checks imported facts, substitutions, clause well-formedness, and the resolution/SAT trace - it does not trust a solver's exit code, so unsoundness in search cannot become unsoundness in accepted results;
\item a proof accepted under one dependency slice cannot silently migrate to another: the hashes pin it down.
\end{itemize}
\note{
\smallskip
\textbf{Presenter script:}
\smallskip
\textbf{Read aloud:}
\begin{itemize}
\setlength{\itemsep}{1pt}
\setlength{\parsep}{0pt}
\setlength{\parskip}{0pt}
\item accepted evidence is normalized into replayable certificate data;
\item the kernel checks imported facts, substitutions, clause well-formedness, and the resolution/SAT trace - it does not trust a solver's exit code, so unsoundness in search cannot become unsoundness in accepted results;
\item a proof accepted under one dependency slice cannot silently migrate to another: the hashes pin it down.
\end{itemize}
After the example, say which design obligation the syntax makes visible.

}
\end{frame}

\begin{frame}[fragile,t]{5.4 - The Same Boundary Tames AI}
\smallskip
\textbf{Legacy Mizar (exact MML excerpt):}
\begin{Verbatim}[fontsize=\small,frame=single,framesep=0.8mm,framerule=0.35pt,rulecolor=\color{black!45},xleftmargin=0.5mm,xrightmargin=0.5mm]
theorem
  for F being non degenerated ZeroOneStr holds 1.F in NonZero F
proof
  let F be non degenerated ZeroOneStr;
  not 1.F in {0.F} by TARSKI:def 1;
  hence thesis by XBOOLE_0:def 5;
end;
\end{Verbatim}
\smallskip
\textbf{Message:}
\begin{itemize}
\setlength{\itemsep}{1pt}
\setlength{\parsep}{0pt}
\setlength{\parskip}{0pt}
\item Citation repair - proposing a missing or sharper \texttt{by} reference - is the canonical safe AI edit: source-local, meaning-preserving, and checked by the verifier like any human edit.
\end{itemize}
\note{
\smallskip
\textbf{Presenter script:}
\smallskip
\textbf{Say:}
\begin{itemize}
\setlength{\itemsep}{1pt}
\setlength{\parsep}{0pt}
\setlength{\parskip}{0pt}
\item Source: current MML \texttt{struct\_0.miz}, lines 637-643.
\item The agent proposes; the verifier and kernel decide. The assistant's strength never enters the trusted base.
\end{itemize}
\smallskip
\textbf{Read aloud:}
\begin{itemize}
\setlength{\itemsep}{1pt}
\setlength{\parsep}{0pt}
\setlength{\parskip}{0pt}
\item Citation repair - proposing a missing or sharper by reference - is the canonical safe AI edit: source-local, meaning-preserving, and checked by the verifier like any human edit.
\end{itemize}
}
\end{frame}

\begin{frame}[fragile,t]{5.5 - Edit Classes: Green, Yellow, Red}
\begin{footnotesize}
\setlength{\tabcolsep}{1.5pt}
\renewcommand{\arraystretch}{1.0}
\begin{tabularx}{\textwidth}{|>{\raggedright\arraybackslash}X|>{\raggedright\arraybackslash}X|>{\raggedright\arraybackslash}X|}
\hline
Class & Examples & Policy \\
\hline
Green & add citation, insert \texttt{qua}, info annotation & auto-proposable, still verified \\
\hline
Yellow & add import, local lemma, registration & proposed with human review \\
\hline
Red & weaken theorem, change definition, add axiom & forbidden to ordinary agents \\
\hline
\end{tabularx}
\end{footnotesize}

\smallskip
\textbf{Forbidden repair (sketch):}
\begin{Verbatim}[fontsize=\small,frame=single,framesep=0.8mm,framerule=0.35pt,rulecolor=\color{black!45},xleftmargin=0.5mm,xrightmargin=0.5mm]
theorem
  for x be Nat holds x + 0 = x or x = 0;
\end{Verbatim}
\smallskip
\textbf{Message:}
\begin{itemize}
\setlength{\itemsep}{1pt}
\setlength{\parsep}{0pt}
\setlength{\parskip}{0pt}
\item Weakening a statement to make its proof easier is a Red edit; at most an agent may flag it for explicit human unsafe-edit review.
\end{itemize}
\note{
\smallskip
\textbf{Presenter script:}
\smallskip
\textbf{Read aloud:}
\begin{itemize}
\setlength{\itemsep}{1pt}
\setlength{\parsep}{0pt}
\setlength{\parskip}{0pt}
\item Weakening a statement to make its proof easier is a Red edit; at most an agent may flag it for explicit human unsafe-edit review.
\end{itemize}
After the example, say which design obligation the syntax makes visible.

For the table, read the row labels first, then the design reason.

}
\end{frame}

\begin{frame}[fragile,t]{5.6 - What Is Preserved, What We Ask}
\smallskip
\textbf{Preserved:}
\begin{itemize}
\setlength{\itemsep}{1pt}
\setlength{\parsep}{0pt}
\setlength{\parskip}{0pt}
\item the de Bruijn discipline: a small checker judges everything;
\item proof text stays declarative and readable - automation maintains the argument, it does not replace it.
\end{itemize}
\smallskip
\textbf{Questions for Bialystok:}
\begin{itemize}
\setlength{\itemsep}{1pt}
\setlength{\parsep}{0pt}
\setlength{\parskip}{0pt}
\item Is the SAT/resolution certificate story convincing for Mizar-style obligations, including clusters and definitional expansions?
\item Which evidence format would the team be most willing to audit?
\end{itemize}
\note{
\smallskip
\textbf{Presenter script:}
\smallskip
\textbf{Read aloud:}
\begin{itemize}
\setlength{\itemsep}{1pt}
\setlength{\parsep}{0pt}
\setlength{\parskip}{0pt}
\item the de Bruijn discipline: a small checker judges everything;
\item proof text stays declarative and readable - automation maintains the argument, it does not replace it.
\item Is the SAT/resolution certificate story convincing for Mizar-style obligations, including clusters and definitional expansions?
\item Which evidence format would the team be most willing to audit?
\end{itemize}
}
\end{frame}

\section{Part 6. Story 5: Verification That Scales}
\begin{frame}[fragile,t]{6.1 - The Pain}
\smallskip
\textbf{Key Points:}
\begin{itemize}
\setlength{\itemsep}{1pt}
\setlength{\parsep}{0pt}
\setlength{\parskip}{0pt}
\item Verifying the whole MML is a batch operation measured in hours.
\item The reuse boundary is the accepted article: a small change re-verifies more than it should.
\item Memory follows the article environment, not the actually used interface.
\item None of this is a defect of the current design - it is what article-level granularity implies once the library is large.
\end{itemize}
\note{
\smallskip
\textbf{Presenter script:}
\smallskip
\textbf{Read aloud:}
\begin{itemize}
\setlength{\itemsep}{1pt}
\setlength{\parsep}{0pt}
\setlength{\parskip}{0pt}
\item Verifying the whole MML is a batch operation measured in hours.
\item The reuse boundary is the accepted article: a small change re-verifies more than it should.
\item Memory follows the article environment, not the actually used interface.
\item None of this is a defect of the current design - it is what article-level granularity implies once the library is large.
\end{itemize}
}
\end{frame}

\begin{frame}[fragile,t]{6.2 - The Evo Answer: Fingerprints And Incrementality}
\smallskip
\textbf{Key Points:}
\begin{itemize}
\setlength{\itemsep}{1pt}
\setlength{\parsep}{0pt}
\setlength{\parskip}{0pt}
\item hash the source, the dependency interfaces, the generated obligations, and the proof witnesses;
\item reuse a result only when fingerprints and verifier policy match;
\item a proof-body-only change does not rebuild importers, because public statements and statuses are unchanged;
\item independent modules, obligations, ATP runs, and kernel checks run in parallel; results are published in canonical order.
\end{itemize}
\smallskip
\textbf{Rule:}
\begin{Verbatim}[fontsize=\small,frame=single,framesep=0.8mm,framerule=0.35pt,rulecolor=\color{black!45},xleftmargin=0.5mm,xrightmargin=0.5mm]
Cache reuse is never proof authority.
A clean build must always be able to reproduce every acceptance.
\end{Verbatim}
\note{
\smallskip
\textbf{Presenter script:}
\smallskip
\textbf{Read aloud:}
\begin{itemize}
\setlength{\itemsep}{1pt}
\setlength{\parsep}{0pt}
\setlength{\parskip}{0pt}
\item hash the source, the dependency interfaces, the generated obligations, and the proof witnesses;
\item reuse a result only when fingerprints and verifier policy match;
\item a proof-body-only change does not rebuild importers, because public statements and statuses are unchanged;
\end{itemize}
After the example, say which design obligation the syntax makes visible.

}
\end{frame}

\begin{frame}[fragile,t]{6.3 - The Evo Answer: A Memory Contract}
\begin{Verbatim}[fontsize=\footnotesize,frame=single,framesep=0.8mm,framerule=0.35pt,rulecolor=\color{black!45},xleftmargin=0.5mm,xrightmargin=0.5mm]
resident memory should scale with:
  active source
  imported public interfaces
  import-filtered indexes
  active module obligations

not with:
  imported proof bodies
  private lemmas outside the interface
  registration data outside the import closure
\end{Verbatim}
\smallskip
\textbf{Message:}
\begin{itemize}
\setlength{\itemsep}{1pt}
\setlength{\parsep}{0pt}
\setlength{\parskip}{0pt}
\item Interfaces are loaded; proof bodies are not. This is what makes whole-MML editing sessions feasible on ordinary hardware.
\end{itemize}
\note{
\smallskip
\textbf{Presenter script:}
\smallskip
\textbf{Read aloud:}
\begin{itemize}
\setlength{\itemsep}{1pt}
\setlength{\parsep}{0pt}
\setlength{\parskip}{0pt}
\item Interfaces are loaded; proof bodies are not. This is what makes whole-MML editing sessions feasible on ordinary hardware.
\end{itemize}
After the example, say which design obligation the syntax makes visible.

}
\end{frame}

\begin{frame}[fragile,t]{6.4 - What Is Preserved, What We Ask}
\smallskip
\textbf{Preserved:}
\begin{itemize}
\setlength{\itemsep}{1pt}
\setlength{\parsep}{0pt}
\setlength{\parskip}{0pt}
\item clean-build semantics: caching and parallelism change speed, never truth;
\item article-style review: what a human reads is still complete source text.
\end{itemize}
\smallskip
\textbf{Questions for Bialystok:}
\begin{itemize}
\setlength{\itemsep}{1pt}
\setlength{\parsep}{0pt}
\setlength{\parskip}{0pt}
\item What clean-build equivalence tests would make incremental verification trustworthy to the team that maintains MML today?
\item Which current maintenance operations (revisions, renamings) should we benchmark for incremental cost?
\end{itemize}
\note{
\smallskip
\textbf{Presenter script:}
\smallskip
\textbf{Read aloud:}
\begin{itemize}
\setlength{\itemsep}{1pt}
\setlength{\parsep}{0pt}
\setlength{\parskip}{0pt}
\item clean-build semantics: caching and parallelism change speed, never truth;
\item article-style review: what a human reads is still complete source text.
\item What clean-build equivalence tests would make incremental verification trustworthy to the team that maintains MML today?
\item Which current maintenance operations (revisions, renamings) should we benchmark for incremental cost?
\end{itemize}
}
\end{frame}

\section{Part 7. Story 6: Templates For Generic Mathematics}
\begin{frame}[fragile,t]{7.1 - The Pain, Part One: Schemes Are Fenced Off}
\smallskip
\textbf{Legacy Mizar (exact MML excerpt):}
\begin{Verbatim}[fontsize=\small,frame=single,framesep=0.8mm,framerule=0.35pt,rulecolor=\color{black!45},xleftmargin=0.5mm,xrightmargin=0.5mm]
scheme
  NatInd { P[Nat] } : for k being Nat holds P[k]
provided
A1: P[0] and
A2: for k be Nat st P[k] holds P[k + 1]
\end{Verbatim}
\smallskip
\textbf{Key Points:}
\begin{itemize}
\setlength{\itemsep}{1pt}
\setlength{\parsep}{0pt}
\setlength{\parskip}{0pt}
\item schemes carry second-order patterns (induction, separation, replacement) and they work - but they are a separate mechanism with separate rules;
\item schemes can parameterize theorems, but not structures, modes, or functors.
\end{itemize}
\note{
\smallskip
\textbf{Presenter script:}
\smallskip
\textbf{Say:}
\begin{itemize}
\setlength{\itemsep}{1pt}
\setlength{\parsep}{0pt}
\setlength{\parskip}{0pt}
\item Source: current MML \texttt{nat\_1.miz}, near line 90 (checked July 2, 2026).
\end{itemize}
\smallskip
\textbf{Read aloud:}
\begin{itemize}
\setlength{\itemsep}{1pt}
\setlength{\parsep}{0pt}
\setlength{\parskip}{0pt}
\item schemes carry second-order patterns (induction, separation, replacement) and they work - but they are a separate mechanism with separate rules;
\item schemes can parameterize theorems, but not structures, modes, or functors.
\end{itemize}
}
\end{frame}

\begin{frame}[fragile,t]{7.2 - The Pain, Part Two: Copy-Paste Algebra}
\smallskip
\textbf{Key Points:}
\begin{itemize}
\setlength{\itemsep}{1pt}
\setlength{\parsep}{0pt}
\setlength{\parskip}{0pt}
\item \texttt{addMagma} and \texttt{multMagma} are the same mathematics twice, related only by naming convention (story 2 met this already);
\item polynomial rings, vector spaces, and matrix theories are re-spelled per carrier because there is no parameterized construction;
\item a theorem proved for one commutative operation is re-proved for \texttt{+} and \texttt{*} separately.
\end{itemize}
\smallskip
\textbf{Message:}
\begin{itemize}
\setlength{\itemsep}{1pt}
\setlength{\parsep}{0pt}
\setlength{\parskip}{0pt}
\item The library pays for the missing generics mechanism in duplicated articles, and every duplicate is a maintenance obligation.
\end{itemize}
\note{
\smallskip
\textbf{Presenter script:}
\smallskip
\textbf{Read aloud:}
\begin{itemize}
\setlength{\itemsep}{1pt}
\setlength{\parsep}{0pt}
\setlength{\parskip}{0pt}
\item addMagma and multMagma are the same mathematics twice, related only by naming convention (story 2 met this already);
\item polynomial rings, vector spaces, and matrix theories are re-spelled per carrier because there is no parameterized construction;
\item a theorem proved for one commutative operation is re-proved for + and * separately.
\item The library pays for the missing generics mechanism in duplicated articles, and every duplicate is a maintenance obligation.
\end{itemize}
}
\end{frame}

\begin{frame}[fragile,t]{7.3 - The Evo Answer: Templates}
\smallskip
\textbf{Mizar Evo (specification example):}
\begin{Verbatim}[fontsize=\small,frame=single,framesep=0.8mm,framerule=0.35pt,rulecolor=\color{black!45},xleftmargin=0.5mm,xrightmargin=0.5mm]
definition
  let T be type;
  struct MagmaStr[T] where
    field carrier -> T;
    field binop -> BinOp of T;
  end;
end;
\end{Verbatim}
\smallskip
\textbf{Key Points:}
\begin{itemize}
\setlength{\itemsep}{1pt}
\setlength{\parsep}{0pt}
\setlength{\parskip}{0pt}
\item a template is an ordinary \texttt{definition} block whose leading \texttt{let} binds parameters: types, values, predicates, or functors;
\item one mechanism covers structures, modes, functors, predicates, theorems, registrations, and algorithms;
\item readable shorthands survive: \texttt{Module over R} is an automatic synonym for \texttt{Module[R]}, \texttt{Subset of X} for \texttt{Subset[X]}.
\end{itemize}
\note{
\smallskip
\textbf{Presenter script:}
\smallskip
\textbf{Read aloud:}
\begin{itemize}
\setlength{\itemsep}{1pt}
\setlength{\parsep}{0pt}
\setlength{\parskip}{0pt}
\item a template is an ordinary definition block whose leading let binds parameters: types, values, predicates, or functors;
\item one mechanism covers structures, modes, functors, predicates, theorems, registrations, and algorithms;
\item readable shorthands survive: Module over R is an automatic synonym for Module[R], Subset of X for Subset[X].
\end{itemize}
After the example, say which design obligation the syntax makes visible.

}
\end{frame}

\begin{frame}[fragile,t]{7.4 - Bounded Parameters And Generic Theorems}
\smallskip
\textbf{Mizar Evo (specification example):}
\begin{Verbatim}[fontsize=\footnotesize,frame=single,framesep=0.8mm,framerule=0.35pt,rulecolor=\color{black!45},xleftmargin=0.5mm,xrightmargin=0.5mm]
definition
  let T be type extends commutative Magma;
  theorem PermProduct[T]:
    for s being FinSequence of T,
        p being Permutation of dom s
    holds Product[T](s) = Product[T](s * p)
  proof ... end;
end;
\end{Verbatim}
\smallskip
\textbf{Message:}
\begin{itemize}
\setlength{\itemsep}{1pt}
\setlength{\parsep}{0pt}
\setlength{\parskip}{0pt}
\item \texttt{type extends commutative Magma} states what the parameter must provide before any proof search begins;
\item the theorem is proved once, for every commutative operation.
\end{itemize}
\note{
\smallskip
\textbf{Presenter script:}
\smallskip
\textbf{Read aloud:}
\begin{itemize}
\setlength{\itemsep}{1pt}
\setlength{\parsep}{0pt}
\setlength{\parskip}{0pt}
\item type extends commutative Magma states what the parameter must provide before any proof search begins;
\item the theorem is proved once, for every commutative operation.
\end{itemize}
After the example, say which design obligation the syntax makes visible.

}
\end{frame}

\begin{frame}[fragile,t]{7.5 - One Proof, Many Instantiations}
\smallskip
\textbf{Instantiation (specification example):}
\begin{Verbatim}[fontsize=\small,frame=single,framesep=0.8mm,framerule=0.35pt,rulecolor=\color{black!45},xleftmargin=0.5mm,xrightmargin=0.5mm]
PermProduct[AddMagma]              :: additive form
PermProduct[commutative MulMagma]  :: multiplicative form

let R be commutative Ring;
PermProduct[R qua AddMagma]        :: R's additive view
PermProduct[R qua MulMagma]        :: R's multiplicative view
\end{Verbatim}
\smallskip
\textbf{Message:}
\begin{itemize}
\setlength{\itemsep}{1pt}
\setlength{\parsep}{0pt}
\setlength{\parskip}{0pt}
\item \texttt{qua} selects the intended view when a ring reaches Magma along two paths - and notation follows the view: the generic \texttt{*} displays as \texttt{+} under the additive instantiation.
\end{itemize}
\note{
\smallskip
\textbf{Presenter script:}
\smallskip
\textbf{Read aloud:}
\begin{itemize}
\setlength{\itemsep}{1pt}
\setlength{\parsep}{0pt}
\setlength{\parskip}{0pt}
\item qua selects the intended view when a ring reaches Magma along two paths - and notation follows the view: the generic * displays as + under the additive instantiation.
\end{itemize}
After the example, say which design obligation the syntax makes visible.

}
\end{frame}

\begin{frame}[fragile,t]{7.6 - Schemes Become Ordinary Templates}
\smallskip
\textbf{Mizar Evo (specification example):}
\begin{Verbatim}[fontsize=\small,frame=single,framesep=0.8mm,framerule=0.35pt,rulecolor=\color{black!45},xleftmargin=0.5mm,xrightmargin=0.5mm]
definition
  let P be pred(Nat);
  theorem NatInduction[P]:
    P(0) & (for n being Nat st P(n) holds P(n+1))
    implies for n being Nat holds P(n)
  proof ... end;
end;
\end{Verbatim}
\smallskip
\textbf{Key Points:}
\begin{itemize}
\setlength{\itemsep}{1pt}
\setlength{\parsep}{0pt}
\setlength{\parskip}{0pt}
\item predicate parameters follow the familiar \texttt{defpred} convention;
\item legacy schemes have a direct, mechanical migration target;
\item instantiation is explicit bracket syntax, so tools and artifacts see exactly which instance a proof uses.
\end{itemize}
\note{
\smallskip
\textbf{Presenter script:}
\smallskip
\textbf{Read aloud:}
\begin{itemize}
\setlength{\itemsep}{1pt}
\setlength{\parsep}{0pt}
\setlength{\parskip}{0pt}
\item predicate parameters follow the familiar defpred convention;
\item legacy schemes have a direct, mechanical migration target;
\item instantiation is explicit bracket syntax, so tools and artifacts see exactly which instance a proof uses.
\end{itemize}
After the example, say which design obligation the syntax makes visible.

}
\end{frame}

\begin{frame}[fragile,t]{7.7 - What Is Preserved, What We Ask}
\smallskip
\textbf{Preserved:}
\begin{itemize}
\setlength{\itemsep}{1pt}
\setlength{\parsep}{0pt}
\setlength{\parskip}{0pt}
\item scheme-style reasoning survives unchanged in power;
\item \texttt{of} / \texttt{over} phrasing keeps mathematical prose readable;
\item first-order discipline: templates are checked instantiation, not a new logic.
\end{itemize}
\smallskip
\textbf{Questions for Bialystok:}
\begin{itemize}
\setlength{\itemsep}{1pt}
\setlength{\parsep}{0pt}
\setlength{\parskip}{0pt}
\item Which MML schemes should be the first migration targets?
\item Are brackets acceptable as the canonical identity form, with \texttt{of}/\texttt{over} as display forms?
\item Where should parameter inference stop and demand explicit \texttt{[T]} or \texttt{qua}?
\end{itemize}
\note{
\smallskip
\textbf{Presenter script:}
\smallskip
\textbf{Read aloud:}
\begin{itemize}
\setlength{\itemsep}{1pt}
\setlength{\parsep}{0pt}
\setlength{\parskip}{0pt}
\item scheme-style reasoning survives unchanged in power;
\item of / over phrasing keeps mathematical prose readable;
\item first-order discipline: templates are checked instantiation, not a new logic.
\item Which MML schemes should be the first migration targets?
\item Are brackets acceptable as the canonical identity form, with of/over as display forms?
\end{itemize}
}
\end{frame}

\section{Part 8. Story 7: Verified Computation With Algorithms}
\begin{frame}[fragile,t]{8.1 - The Pain}
\smallskip
\textbf{Key Points:}
\begin{itemize}
\setlength{\itemsep}{1pt}
\setlength{\parsep}{0pt}
\setlength{\parskip}{0pt}
\item Current Mizar can define \texttt{Gcd} and prove its theory - but cannot compute \texttt{gcd(48, 18)}; every numeric fact needs a hand-written proof chain.
\item There is no checked connection between MML mathematics and executable code; verified-algorithm work must leave the system entirely.
\item This is a boundary of the design, not a defect: Mizar chose to be a proof language. The question is whether that boundary still serves us.
\end{itemize}
\smallskip
\textbf{Key Phrase:}
\begin{Verbatim}[fontsize=\small,frame=single,framesep=0.8mm,framerule=0.35pt,rulecolor=\color{black!45},xleftmargin=0.5mm,xrightmargin=0.5mm]
The library describes computation.
It cannot perform or export it.
\end{Verbatim}
\note{
\smallskip
\textbf{Presenter script:}
\smallskip
\textbf{Read aloud:}
\begin{itemize}
\setlength{\itemsep}{1pt}
\setlength{\parsep}{0pt}
\setlength{\parskip}{0pt}
\item Current Mizar can define Gcd and prove its theory - but cannot compute gcd(48, 18); every numeric fact needs a hand-written proof chain.
\item There is no checked connection between MML mathematics and executable code; verified-algorithm work must leave the system entirely.
\item This is a boundary of the design, not a defect: Mizar chose to be a proof language. The question is whether that boundary still serves us.
\end{itemize}
After the example, say which design obligation the syntax makes visible.

}
\end{frame}

\begin{frame}[fragile,t]{8.2 - The Evo Answer: Algorithms With Contracts}
\smallskip
\textbf{Mizar Evo (specification example, condensed from spec section 20.12):}
\begin{Verbatim}[fontsize=\footnotesize,frame=single,framesep=0.8mm,framerule=0.35pt,rulecolor=\color{black!45},xleftmargin=0.5mm,xrightmargin=0.5mm]
definition
  let a, b be Nat;
  terminating algorithm euclid_gcd(a, b) -> Nat
    requires a >= 1 & b >= 1
    ensures result = Gcd(a, b)
  do
    var x := a;  var y := b;
    while y <> 0 do
      invariant x >= 1 & y >= 0 & Gcd(a, b) = Gcd(x, y);
      decreasing y;
      const r := x mod y;  x := y;  y := r;
    end;
    return x;
  end;
end;
\end{Verbatim}
\smallskip
\textbf{Message:}
\begin{itemize}
\setlength{\itemsep}{1pt}
\setlength{\parsep}{0pt}
\setlength{\parskip}{0pt}
\item \texttt{ensures} and the invariant cite the mathematical \texttt{Gcd}: no circularity.
\end{itemize}
\note{
\smallskip
\textbf{Presenter script:}
\smallskip
\textbf{Say:}
\begin{itemize}
\setlength{\itemsep}{1pt}
\setlength{\parsep}{0pt}
\setlength{\parskip}{0pt}
\item The contract speaks the mathematical language: the algorithm computes, while the library functor \texttt{Gcd} specifies. Termination comes from the \texttt{decreasing} measure on \texttt{y}.
\end{itemize}
\smallskip
\textbf{Read aloud:}
\begin{itemize}
\setlength{\itemsep}{1pt}
\setlength{\parsep}{0pt}
\setlength{\parskip}{0pt}
\item ensures and the invariant cite the mathematical Gcd: no circularity.
\end{itemize}
}
\end{frame}

\begin{frame}[fragile,t]{8.3 - The Evo Answer: Proof By Computation}
\smallskip
\textbf{Mizar Evo (specification example):}
\begin{Verbatim}[fontsize=\small,frame=single,framesep=0.8mm,framerule=0.35pt,rulecolor=\color{black!45},xleftmargin=0.5mm,xrightmargin=0.5mm]
theorem EuclidGcd12_8:  euclid_gcd(12, 8)  = 4  by computation;
theorem EuclidGcd100_75: euclid_gcd(100, 75) = 25 by computation;

theorem Fact10: factorial(10) = 3628800
proof
  thus thesis by computation(steps: 100000);
end;
\end{Verbatim}
\smallskip
\textbf{Key Points:}
\begin{itemize}
\setlength{\itemsep}{1pt}
\setlength{\parsep}{0pt}
\setlength{\parskip}{0pt}
\item the Mizar Virtual Machine (MVM) evaluates ground goals during verification, under explicit step, time, and depth budgets;
\item only ground equalities and ground predicates qualify; everything else still requires a classical proof;
\item imported algorithms are opaque: downstream proofs may use only their \texttt{ensures} contract, never their body.
\end{itemize}
\note{
\smallskip
\textbf{Presenter script:}
\smallskip
\textbf{Read aloud:}
\begin{itemize}
\setlength{\itemsep}{1pt}
\setlength{\parsep}{0pt}
\setlength{\parskip}{0pt}
\item the Mizar Virtual Machine (MVM) evaluates ground goals during verification, under explicit step, time, and depth budgets;
\item only ground equalities and ground predicates qualify; everything else still requires a classical proof;
\item imported algorithms are opaque: downstream proofs may use only their ensures contract, never their body.
\end{itemize}
After the example, say which design obligation the syntax makes visible.

}
\end{frame}

\begin{frame}[fragile,t]{8.4 - The Evo Answer: Termination Buys Recursion}
\smallskip
\textbf{Mizar Evo (specification example):}
\begin{Verbatim}[fontsize=\footnotesize,frame=single,framesep=0.8mm,framerule=0.35pt,rulecolor=\color{black!45},xleftmargin=0.5mm,xrightmargin=0.5mm]
definition
  let n be Nat;
  terminating algorithm factorial(n) -> Nat
    decreasing n
  do
    if n = 0 do return 1; end;
    return n * factorial(n - 1);
  end;
end;
\end{Verbatim}
\smallskip
\textbf{Key Points:}
\begin{itemize}
\setlength{\itemsep}{1pt}
\setlength{\parsep}{0pt}
\setlength{\parskip}{0pt}
\item ordinary \texttt{func} definitions are first-order definitional extensions and cannot be recursive;
\item a \texttt{terminating} algorithm - once its termination obligations are discharged - is promoted to a genuine functor, usable in any proof;
\item this is the only door through which recursion enters the mathematical layer, and it is a proof-shaped door.
\end{itemize}
\note{
\smallskip
\textbf{Presenter script:}
\smallskip
\textbf{Read aloud:}
\begin{itemize}
\setlength{\itemsep}{1pt}
\setlength{\parsep}{0pt}
\setlength{\parskip}{0pt}
\item ordinary func definitions are first-order definitional extensions and cannot be recursive;
\item a terminating algorithm - once its termination obligations are discharged - is promoted to a genuine functor, usable in any proof;
\item this is the only door through which recursion enters the mathematical layer, and it is a proof-shaped door.
\end{itemize}
After the example, say which design obligation the syntax makes visible.

}
\end{frame}

\begin{frame}[fragile,t]{8.5 - Computation Never Redefines Truth}
\smallskip
\textbf{Key Points:}
\begin{itemize}
\setlength{\itemsep}{1pt}
\setlength{\parsep}{0pt}
\setlength{\parskip}{0pt}
\item contracts, invariants, and termination measures generate verification conditions that pass through the same ATP-plus-kernel boundary as ordinary theorems (story 4);
\item \texttt{by computation} is MVM replay under budgets, not solver trust;
\item code extraction (to runtime targets) is strictly downstream of verified artifacts and can never feed back into acceptance.
\end{itemize}
\smallskip
\textbf{Message:}
\begin{itemize}
\setlength{\itemsep}{1pt}
\setlength{\parsep}{0pt}
\setlength{\parskip}{0pt}
\item Algorithms widen what the library can express; they do not touch what it means for a theorem to be accepted.
\end{itemize}
\note{
\smallskip
\textbf{Presenter script:}
\smallskip
\textbf{Read aloud:}
\begin{itemize}
\setlength{\itemsep}{1pt}
\setlength{\parsep}{0pt}
\setlength{\parskip}{0pt}
\item contracts, invariants, and termination measures generate verification conditions that pass through the same ATP-plus-kernel boundary as ordinary theorems (story 4);
\item by computation is MVM replay under budgets, not solver trust;
\item code extraction (to runtime targets) is strictly downstream of verified artifacts and can never feed back into acceptance.
\item Algorithms widen what the library can express; they do not touch what it means for a theorem to be accepted.
\end{itemize}
}
\end{frame}

\begin{frame}[fragile,t]{8.6 - What Is Preserved, What We Ask}
\smallskip
\textbf{Preserved:}
\begin{itemize}
\setlength{\itemsep}{1pt}
\setlength{\parsep}{0pt}
\setlength{\parskip}{0pt}
\item the theorem language is unchanged; algorithms live in \texttt{definition} blocks and interact with proofs only through contracts and verified promotion;
\item first-order set-theoretic foundations stay intact.
\end{itemize}
\smallskip
\textbf{Questions for Bialystok:}
\begin{itemize}
\setlength{\itemsep}{1pt}
\setlength{\parsep}{0pt}
\setlength{\parskip}{0pt}
\item Which computational examples would demonstrate value to mathematicians without shifting the culture toward programming?
\item Are there MML areas (number theory, combinatorics, finite structures) where \texttt{by computation} would immediately shorten real proofs?
\item Which extraction targets matter first, if any?
\end{itemize}
\note{
\smallskip
\textbf{Presenter script:}
\smallskip
\textbf{Read aloud:}
\begin{itemize}
\setlength{\itemsep}{1pt}
\setlength{\parsep}{0pt}
\setlength{\parskip}{0pt}
\item the theorem language is unchanged; algorithms live in definition blocks and interact with proofs only through contracts and verified promotion;
\item first-order set-theoretic foundations stay intact.
\item Which computational examples would demonstrate value to mathematicians without shifting the culture toward programming?
\item Are there MML areas (number theory, combinatorics, finite structures) where by computation would immediately shorten real proofs?
\item Which extraction targets matter first, if any?
\end{itemize}
}
\end{frame}

\section{Part 9. Story 8: A Library You Can Cite}
\begin{frame}[fragile,t]{9.1 - The Pain}
\smallskip
\textbf{Key Points:}
\begin{itemize}
\setlength{\itemsep}{1pt}
\setlength{\parsep}{0pt}
\setlength{\parskip}{0pt}
\item Formalized Mathematics gives Mizar something rare: a scholarly, citable publication layer over a formal library.
\item But article identity and library organization are tightly coupled: refactoring the library strains published article structure, and package reuse has no journal-facing identity at all.
\item Exposition wants narrative order; reuse wants dependency order. One structure cannot optimize both.
\end{itemize}
\note{
\smallskip
\textbf{Presenter script:}
\smallskip
\textbf{Read aloud:}
\begin{itemize}
\setlength{\itemsep}{1pt}
\setlength{\parsep}{0pt}
\setlength{\parskip}{0pt}
\item Formalized Mathematics gives Mizar something rare: a scholarly, citable publication layer over a formal library.
\item But article identity and library organization are tightly coupled: refactoring the library strains published article structure, and package reuse has no journal-facing identity at all.
\item Exposition wants narrative order; reuse wants dependency order. One structure cannot optimize both.
\end{itemize}
}
\end{frame}

\begin{frame}[fragile,t]{9.2 - The Evo Answer: Linked, Not Merged}
\begin{Verbatim}[fontsize=\small,frame=single,framesep=0.8mm,framerule=0.35pt,rulecolor=\color{black!45},xleftmargin=0.5mm,xrightmargin=0.5mm]
Formalized Mathematics theorem
  -> library package (name, version)
  -> module path and theorem FQN
  -> statement fingerprint
  -> verified artifact hash
  -> origin id (e.g. mizar:GROUP_1:...)
\end{Verbatim}
\smallskip
\textbf{Message:}
\begin{itemize}
\setlength{\itemsep}{1pt}
\setlength{\parsep}{0pt}
\setlength{\parskip}{0pt}
\item each layer answers a different question: scholarly citation, current location, semantic drift detection, reproducible verification, and historical continuity with MML and past Formalized Mathematics volumes.
\end{itemize}
\note{
\smallskip
\textbf{Presenter script:}
\smallskip
\textbf{Read aloud:}
\begin{itemize}
\setlength{\itemsep}{1pt}
\setlength{\parsep}{0pt}
\setlength{\parskip}{0pt}
\item each layer answers a different question: scholarly citation, current location, semantic drift detection, reproducible verification, and historical continuity with MML and past Formalized Mathematics volumes.
\end{itemize}
After the example, say which design obligation the syntax makes visible.

}
\end{frame}

\begin{frame}[fragile,t]{9.3 - Who Gains What}
\begin{footnotesize}
\setlength{\tabcolsep}{1.5pt}
\renewcommand{\arraystretch}{1.0}
\begin{tabularx}{\textwidth}{|>{\raggedright\arraybackslash}X|>{\raggedright\arraybackslash}X|}
\hline
Audience & Gain \\
\hline
readers & prose stays primary; formal source is one click away \\
\hline
maintainers & refactoring no longer rewrites published exposition \\
\hline
authors & articles cite stable identities, not file layouts \\
\hline
AI tools & prose for retrieval, fingerprints for exact context \\
\hline
\end{tabularx}
\end{footnotesize}

\note{
\smallskip
\textbf{Presenter script:}
For the table, read the row labels first, then the design reason.

Use this as the transition: Who Gains What. Point to the visible example, question, or diagram before moving on.

}
\end{frame}

\begin{frame}[fragile,t]{9.4 - What Is Preserved, What We Ask}
\smallskip
\textbf{Preserved:}
\begin{itemize}
\setlength{\itemsep}{1pt}
\setlength{\parsep}{0pt}
\setlength{\parskip}{0pt}
\item Formalized Mathematics remains a real journal with review and exposition;
\item origin metadata keeps continuity with every existing MML citation.
\end{itemize}
\smallskip
\textbf{Questions for Bialystok:}
\begin{itemize}
\setlength{\itemsep}{1pt}
\setlength{\parsep}{0pt}
\setlength{\parskip}{0pt}
\item Which identity should be primary in user-facing citations: article label, library FQN, or origin id?
\item How should existing Formalized Mathematics articles link to migrated modules - retroactively, on revision, or not at all?
\end{itemize}
\note{
\smallskip
\textbf{Presenter script:}
\smallskip
\textbf{Read aloud:}
\begin{itemize}
\setlength{\itemsep}{1pt}
\setlength{\parsep}{0pt}
\setlength{\parskip}{0pt}
\item Formalized Mathematics remains a real journal with review and exposition;
\item origin metadata keeps continuity with every existing MML citation.
\item Which identity should be primary in user-facing citations: article label, library FQN, or origin id?
\item How should existing Formalized Mathematics articles link to migrated modules - retroactively, on revision, or not at all?
\end{itemize}
}
\end{frame}

\section{Part 10. Architecture In One Picture}
\begin{frame}[fragile,t]{10.1 - The Pipeline}
\begin{Verbatim}[fontsize=\small,frame=single,framesep=0.8mm,framerule=0.35pt,rulecolor=\color{black!45},xleftmargin=0.5mm,xrightmargin=0.5mm]
Source
  -> TokenStream -> SurfaceAst -> ResolvedAst -> TypedAst
  -> CoreIr -> VcIr -> AtpProblem -> ProofCertificate
  -> VerifiedArtifact
\end{Verbatim}
\smallskip
\textbf{Message:}
\begin{itemize}
\setlength{\itemsep}{1pt}
\setlength{\parsep}{0pt}
\setlength{\parskip}{0pt}
\item every boundary states who owns a fact, which artifact records it, and what must be recomputed when it changes - that is the entire point.
\end{itemize}
\note{
\smallskip
\textbf{Presenter script:}
\smallskip
\textbf{Read aloud:}
\begin{itemize}
\setlength{\itemsep}{1pt}
\setlength{\parsep}{0pt}
\setlength{\parskip}{0pt}
\item every boundary states who owns a fact, which artifact records it, and what must be recomputed when it changes - that is the entire point.
\end{itemize}
After the example, say which design obligation the syntax makes visible.

}
\end{frame}

\begin{frame}[fragile,t]{10.2 - Responsibility Split}
\begin{footnotesize}
\setlength{\tabcolsep}{1.5pt}
\renewcommand{\arraystretch}{1.0}
\begin{tabularx}{\textwidth}{|>{\raggedright\arraybackslash}X|>{\raggedright\arraybackslash}X|}
\hline
Layer & Responsibility \\
\hline
frontend & lexing, parsing, recovery \\
\hline
resolver & imports, names, labels, namespaces \\
\hline
checker & soft types, clusters, registrations, overloads \\
\hline
elaborator & core logical representation \\
\hline
VC generator & proof and algorithm obligations \\
\hline
ATP layer plus kernel & untrusted search, then acceptance by checking \\
\hline
artifact emitter & stable outputs for tools and dependents \\
\hline
\end{tabularx}
\end{footnotesize}

\note{
\smallskip
\textbf{Presenter script:}
For the table, read the row labels first, then the design reason.

Use this as the transition: Responsibility Split. Point to the visible example, question, or diagram before moving on.

}
\end{frame}

\begin{frame}[fragile,t]{10.3 - Where The Eight Stories Live}
\begin{footnotesize}
\setlength{\tabcolsep}{1.5pt}
\renewcommand{\arraystretch}{1.0}
\begin{tabularx}{\textwidth}{|>{\raggedright\arraybackslash}X|>{\raggedright\arraybackslash}X|}
\hline
Story & Pipeline home \\
\hline
dependencies & resolver, package manager \\
\hline
structures and automation & checker (inheritance and cluster graphs, traces) \\
\hline
search vs trust & ATP layer, kernel, certificates \\
\hline
scale & artifacts, fingerprints, scheduler \\
\hline
templates & elaborator (checked instantiation) \\
\hline
algorithms & VC generator, MVM \\
\hline
publication & artifact emitter, doc generation \\
\hline
\end{tabularx}
\end{footnotesize}

\smallskip
\textbf{Message:}
\begin{itemize}
\setlength{\itemsep}{1pt}
\setlength{\parsep}{0pt}
\setlength{\parskip}{0pt}
\item the stories are not eight separate projects; they are one pipeline seen from eight user-visible pains.
\end{itemize}
\note{
\smallskip
\textbf{Presenter script:}
\smallskip
\textbf{Read aloud:}
\begin{itemize}
\setlength{\itemsep}{1pt}
\setlength{\parsep}{0pt}
\setlength{\parskip}{0pt}
\item the stories are not eight separate projects; they are one pipeline seen from eight user-visible pains.
\end{itemize}
For the table, read the row labels first, then the design reason.

}
\end{frame}

\begin{frame}[fragile,t]{10.4 - Testing The Trust Boundary}
\smallskip
\textbf{Key Phrase:}
\begin{Verbatim}[fontsize=\small,frame=single,framesep=0.8mm,framerule=0.35pt,rulecolor=\color{black!45},xleftmargin=0.5mm,xrightmargin=0.5mm]
Reject what must not pass
before
Accept everything that should pass
\end{Verbatim}
\smallskip
\textbf{Key Points:}
\begin{itemize}
\setlength{\itemsep}{1pt}
\setlength{\parsep}{0pt}
\setlength{\parskip}{0pt}
\item soundness bugs outrank parser gaps; kernel-adjacent tests emphasize malformed and failing evidence first;
\item accepted-language coverage grows behind that shield.
\end{itemize}
\note{
\smallskip
\textbf{Presenter script:}
\smallskip
\textbf{Read aloud:}
\begin{itemize}
\setlength{\itemsep}{1pt}
\setlength{\parsep}{0pt}
\setlength{\parskip}{0pt}
\item soundness bugs outrank parser gaps; kernel-adjacent tests emphasize malformed and failing evidence first;
\item accepted-language coverage grows behind that shield.
\end{itemize}
After the example, say which design obligation the syntax makes visible.

}
\end{frame}

\section{Part 11. Roadmap And Collaboration}
\begin{frame}[fragile,t]{11.1 - Migration Is A Research Program}
\smallskip
\textbf{Phases:}
\begin{enumerate}
\setlength{\itemsep}{1pt}
\setlength{\parsep}{0pt}
\setlength{\parskip}{0pt}
\item End of 2026, alpha: frontend and parser for a core subset, import and module resolution prototype, structured diagnostics, early artifacts.
\item 2027, migration laboratory: 3-5 representative MML articles translated by hand and by script; every mismatch recorded as a classified issue.
\item 2027-2028, expansion: foundational set and relation fragments, functions and binary operations, algebraic structures, then dependency cones around successful fragments.
\end{enumerate}
\smallskip
\textbf{Non-goals for the alpha:}
\begin{itemize}
\setlength{\itemsep}{1pt}
\setlength{\parsep}{0pt}
\setlength{\parskip}{0pt}
\item full MML verification, final compatibility layer, stable AI protocol.
\end{itemize}
\note{
\smallskip
\textbf{Presenter script:}
\smallskip
\textbf{Read aloud:}
\begin{itemize}
\setlength{\itemsep}{1pt}
\setlength{\parsep}{0pt}
\setlength{\parskip}{0pt}
\item End of 2026, alpha: frontend and parser for a core subset, import and module resolution prototype, structured diagnostics, early artifacts.
\item 2027, migration laboratory: 3-5 representative MML articles translated by hand and by script; every mismatch recorded as a classified issue.
\item 2027-2028, expansion: foundational set and relation fragments, functions and binary operations, algebraic structures, then dependency cones around successful fragments.
\item full MML verification, final compatibility layer, stable AI protocol.
\end{itemize}
}
\end{frame}

\begin{frame}[fragile,t]{11.2 - What We Will Measure}
\smallskip
\textbf{Key Points:}
\begin{itemize}
\setlength{\itemsep}{1pt}
\setlength{\parsep}{0pt}
\setlength{\parskip}{0pt}
\item translated articles and lines; accepted parser subset;
\item resolved imports versus unresolved dependencies;
\item obligations closed deterministically versus by ATP-plus-certificate;
\item memory and wall-clock per module, incremental versus clean;
\item compatibility decisions that required human judgment.
\end{itemize}
\note{
\smallskip
\textbf{Presenter script:}
\smallskip
\textbf{Read aloud:}
\begin{itemize}
\setlength{\itemsep}{1pt}
\setlength{\parsep}{0pt}
\setlength{\parskip}{0pt}
\item translated articles and lines; accepted parser subset;
\item resolved imports versus unresolved dependencies;
\item obligations closed deterministically versus by ATP-plus-certificate;
\item memory and wall-clock per module, incremental versus clean;
\item compatibility decisions that required human judgment.
\end{itemize}
}
\end{frame}

\begin{frame}[fragile,t]{11.3 - Compatibility Policy And Risks}
\smallskip
\textbf{Policy:}
\begin{itemize}
\setlength{\itemsep}{1pt}
\setlength{\parsep}{0pt}
\setlength{\parskip}{0pt}
\item preserve theorem identity through origin metadata;
\item keep compatibility aliases where they help migration;
\item record every divergence from old behavior with a reason and a test.
\end{itemize}
\begin{footnotesize}
\setlength{\tabcolsep}{1.5pt}
\renewcommand{\arraystretch}{1.0}
\begin{tabularx}{\textwidth}{|>{\raggedright\arraybackslash}X|>{\raggedright\arraybackslash}X|}
\hline
Risk & Mitigation \\
\hline
compatibility work consumes the project & representative slices first, no big-bang translation \\
\hline
registrations behave differently & trace artifacts and comparison reports early \\
\hline
package layout breaks journal links & origin metadata, article-to-library identifiers \\
\hline
AI edits hide migration mistakes & Red edits stay forbidden; verifier artifacts required \\
\hline
\end{tabularx}
\end{footnotesize}

\note{
\smallskip
\textbf{Presenter script:}
\smallskip
\textbf{Read aloud:}
\begin{itemize}
\setlength{\itemsep}{1pt}
\setlength{\parsep}{0pt}
\setlength{\parskip}{0pt}
\item preserve theorem identity through origin metadata;
\item keep compatibility aliases where they help migration;
\item record every divergence from old behavior with a reason and a test.
\end{itemize}
For the table, read the row labels first, then the design reason.

}
\end{frame}

\begin{frame}[fragile,t]{11.4 - What September 2026 Should Produce}
\smallskip
\textbf{Key Points:}
\begin{itemize}
\setlength{\itemsep}{1pt}
\setlength{\parsep}{0pt}
\setlength{\parskip}{0pt}
\item a prioritized list of migration benchmark articles;
\item agreement on compatibility metadata needs;
\item review notes on the eight stories, especially structures and clusters;
\item a first paper outline;
\item a shared list of objections and risks.
\end{itemize}
\note{
\smallskip
\textbf{Presenter script:}
\smallskip
\textbf{Read aloud:}
\begin{itemize}
\setlength{\itemsep}{1pt}
\setlength{\parsep}{0pt}
\setlength{\parskip}{0pt}
\item a prioritized list of migration benchmark articles;
\item agreement on compatibility metadata needs;
\item review notes on the eight stories, especially structures and clusters;
\item a first paper outline;
\item a shared list of objections and risks.
\end{itemize}
}
\end{frame}

\begin{frame}[fragile,t]{11.5 - What We Ask Of You}
\smallskip
\textbf{Questions:}
\begin{itemize}
\setlength{\itemsep}{1pt}
\setlength{\parsep}{0pt}
\setlength{\parskip}{0pt}
\item Which MML articles are small but structurally representative?
\item Which idioms are culturally essential, beyond technical convenience?
\item Which of the eight stories is most wrong, and why?
\item What migration result would convince the community that Evo is serious?
\end{itemize}
\note{
\smallskip
\textbf{Presenter script:}
\smallskip
\textbf{Read aloud:}
\begin{itemize}
\setlength{\itemsep}{1pt}
\setlength{\parsep}{0pt}
\setlength{\parskip}{0pt}
\item Which MML articles are small but structurally representative?
\item Which idioms are culturally essential, beyond technical convenience?
\item Which of the eight stories is most wrong, and why?
\item What migration result would convince the community that Evo is serious?
\end{itemize}
}
\end{frame}

\section{Part 12. Closing}
\begin{frame}[fragile,t]{12.1 - The Running Question, Again}
\smallskip
\textbf{Slide question:}
\begin{Verbatim}[fontsize=\small,frame=single,framesep=0.8mm,framerule=0.35pt,rulecolor=\color{black!45},xleftmargin=0.5mm,xrightmargin=0.5mm]
What must Mizar Evo preserve so that the Mizar community
still recognizes it as Mizar?
\end{Verbatim}
\note{
\smallskip
\textbf{Presenter script:}
\smallskip
\textbf{Say:}
\begin{itemize}
\setlength{\itemsep}{1pt}
\setlength{\parsep}{0pt}
\setlength{\parskip}{0pt}
\item Return to the opening example: the structure that still reads as Mizar.
\item Invite disagreement story by story, not only in general terms.
\end{itemize}
}
\end{frame}

\begin{frame}[fragile,t]{12.2 - Closing}
\smallskip
\textbf{Final slide:}
\begin{Verbatim}[fontsize=\small,frame=single,framesep=0.8mm,framerule=0.35pt,rulecolor=\color{black!45},xleftmargin=0.5mm,xrightmargin=0.5mm]
Mizar Evo should be modern where scale demands it,
and conservative where Mizar's mathematical identity depends on it.
\end{Verbatim}
\note{
\smallskip
\textbf{Presenter script:}
After the example, say which design obligation the syntax makes visible.

Use this as the transition: Closing. Point to the visible example, question, or diagram before moving on.

}
\end{frame}

\appendix
\section{Backup A. Prepared Exact Examples}
\begin{frame}[fragile,t]{Backup A. Prepared Exact Examples (1/2)}
\smallskip
\textbf{Prepared short excerpts for Beamer conversion:}
\begin{scriptsize}
\setlength{\tabcolsep}{1.5pt}
\renewcommand{\arraystretch}{1.0}
\begin{tabularx}{\textwidth}{|>{\raggedright\arraybackslash}X|>{\raggedright\arraybackslash}X|>{\raggedright\arraybackslash}X|>{\raggedright\arraybackslash}X|}
\hline
Purpose & Source & Lines & Used in \\
\hline
Structure definition & \texttt{algstr\_0.miz} & 37-40 & Frames 0.2, 3.1 \\
\hline
Article environment & \texttt{algstr\_0.miz} & 15-25 & Frame 2.1 \\
\hline
Registration/cluster & \texttt{algstr\_0.miz} & 104-109 & Frame 4.1 \\
\hline
\end{tabularx}
\end{scriptsize}

\smallskip
\textbf{Prepared short excerpts for Beamer conversion:}
\begin{scriptsize}
\setlength{\tabcolsep}{1.5pt}
\renewcommand{\arraystretch}{1.0}
\begin{tabularx}{\textwidth}{|>{\raggedright\arraybackslash}X|>{\raggedright\arraybackslash}X|>{\raggedright\arraybackslash}X|>{\raggedright\arraybackslash}X|}
\hline
Purpose & Source & Lines & Used in \\
\hline
Proof citation & \texttt{struct\_0.miz} & 637-643 & Frame 5.4 \\
\hline
Induction scheme & \texttt{nat\_1.miz} & near 90 & Frame 7.1 \\
\hline
\end{tabularx}
\end{scriptsize}

\smallskip
\textbf{Source URLs:}
\begin{itemize}
\setlength{\itemsep}{1pt}
\setlength{\parsep}{0pt}
\setlength{\parskip}{0pt}
\item \texttt{ALGSTR\_0}: <https://mizar.uwb.edu.pl/version/current/mml/algstr\_0.miz>
\item \texttt{STRUCT\_0}: <https://mizar.uwb.edu.pl/version/current/mml/struct\_0.miz>
\item \texttt{NAT\_1}: <https://mizar.uwb.edu.pl/version/current/mml/nat\_1.miz>
\end{itemize}
\note{
\smallskip
\textbf{Presenter script:}
\smallskip
\textbf{Read aloud:}
\begin{itemize}
\setlength{\itemsep}{1pt}
\setlength{\parsep}{0pt}
\setlength{\parskip}{0pt}
\item ALGSTR\_0: <https://mizar.uwb.edu.pl/version/current/mml/algstr\_0.miz>
\item STRUCT\_0: <https://mizar.uwb.edu.pl/version/current/mml/struct\_0.miz>
\item NAT\_1: <https://mizar.uwb.edu.pl/version/current/mml/nat\_1.miz>
\end{itemize}
For the table, read the row labels first, then the design reason.

}
\end{frame}

\begin{frame}[fragile,t]{Backup A. Prepared Exact Examples (2/2)}
\smallskip
\textbf{Attribution note:}
\begin{itemize}
\setlength{\itemsep}{1pt}
\setlength{\parsep}{0pt}
\setlength{\parskip}{0pt}
\item MML plain-text files state GPL-3.0-or-later / CC-BY-SA-3.0-or-later terms; keep article attribution, URLs, and line numbers in speaker notes.
\end{itemize}
\note{
\smallskip
\textbf{Presenter script:}
\smallskip
\textbf{Read aloud:}
\begin{itemize}
\setlength{\itemsep}{1pt}
\setlength{\parsep}{0pt}
\setlength{\parskip}{0pt}
\item MML plain-text files state GPL-3.0-or-later / CC-BY-SA-3.0-or-later terms; keep article attribution, URLs, and line numbers in speaker notes.
\end{itemize}
}
\end{frame}

\section{Backup B. Specification Reference Map}
\begin{frame}[fragile,t]{Backup B. Specification Reference Map (1/2)}
The slides show examples only. The authoritative grammar and semantics live in the specification; this map replaces the EBNF that earlier drafts put on slides.

\begin{footnotesize}
\setlength{\tabcolsep}{1.5pt}
\renewcommand{\arraystretch}{1.0}
\begin{tabularx}{\textwidth}{|>{\raggedright\arraybackslash}X|>{\raggedright\arraybackslash}X|}
\hline
Topic & Specification source (under \texttt{doc/spec/en/}) \\
\hline
Modules and imports & \texttt{12.modules\_and\_namespaces.md} \\
\hline
Structures and inheritance & \texttt{05.structures.md} \\
\hline
Attributes and modes & \texttt{06.attributes.md}, \texttt{07.modes.md} \\
\hline
Registrations and reductions & \texttt{17.clusters\_and\_registrations.md} \\
\hline
Templates and schemes & \texttt{18.templates.md} \\
\hline
Algorithms and MVM & \texttt{20.algorithm\_and\_verification.md} \\
\hline
Packages and artifacts & \texttt{23.package\_management\_and\_build\_system.md} \\
\hline
\end{tabularx}
\end{footnotesize}

\note{
\smallskip
\textbf{Presenter script:}
\smallskip
\textbf{Read aloud:}
\begin{itemize}
\setlength{\itemsep}{1pt}
\setlength{\parsep}{0pt}
\setlength{\parskip}{0pt}
\item The slides show examples only. The authoritative grammar and semantics live
\item in the specification; this map replaces the EBNF that earlier drafts put on
\item slides.
\end{itemize}
For the table, read the row labels first, then the design reason.

}
\end{frame}

\begin{frame}[fragile,t]{Backup B. Specification Reference Map (2/2)}
\begin{footnotesize}
\setlength{\tabcolsep}{1.5pt}
\renewcommand{\arraystretch}{1.0}
\begin{tabularx}{\textwidth}{|>{\raggedright\arraybackslash}X|>{\raggedright\arraybackslash}X|}
\hline
Topic & Specification source (under \texttt{doc/spec/en/}) \\
\hline
Cross-chapter grammar & \texttt{appendix\_a.grammar\_summary.md} \\
\hline
Worked library sketches & \texttt{sample\_codes.md} \\
\hline
\end{tabularx}
\end{footnotesize}

\note{
\smallskip
\textbf{Presenter script:}
For the table, read the row labels first, then the design reason.

Use this as the transition: Backup B. Specification Reference Map (2/2). Point to the visible example, question, or diagram before moving on.

}
\end{frame}

\section{Backup C. Diagram List}
\begin{frame}[fragile,t]{Backup C. Diagram List}
\smallskip
\textbf{Required diagrams for the final deck:}
\begin{enumerate}
\setlength{\itemsep}{1pt}
\setlength{\parsep}{0pt}
\setlength{\parskip}{0pt}
\item Three pressures on a proven design (Part 1).
\item Environment-to-import migration (story 1).
\item Structure inheritance and diamond coherence (story 2).
\item Reasoning boundary: semantics / ATP search / kernel checking (story 4).
\item Certificate object and replay (story 4).
\item Incremental fingerprint graph (story 5).
\item Formalized Mathematics article-to-library link model (story 8).
\end{enumerate}
\smallskip
\textbf{Required diagrams for the final deck:}
\begin{enumerate}
\setlength{\itemsep}{1pt}
\setlength{\parsep}{0pt}
\setlength{\parskip}{0pt}
\item Full pipeline with story overlay (Part 10).
\item Roadmap timeline (Part 11).
\end{enumerate}
\note{
\smallskip
\textbf{Presenter script:}
\smallskip
\textbf{Read aloud:}
\begin{itemize}
\setlength{\itemsep}{1pt}
\setlength{\parsep}{0pt}
\setlength{\parskip}{0pt}
\item Three pressures on a proven design (Part 1).
\item Environment-to-import migration (story 1).
\item Structure inheritance and diamond coherence (story 2).
\item Reasoning boundary: semantics / ATP search / kernel checking (story 4).
\item Certificate object and replay (story 4).
\end{itemize}
}
\end{frame}

\section{Backup D. Paper Outline Seed}
\begin{frame}[fragile,t]{Backup D. Paper Outline Seed (1/2)}
\smallskip
\textbf{Possible paper title:}
\begin{Verbatim}[fontsize=\small,frame=single,framesep=0.8mm,framerule=0.35pt,rulecolor=\color{black!45},xleftmargin=0.5mm,xrightmargin=0.5mm]
Mizar Evo: Readable, AI-Ready, and Scalable Formal Mathematics
\end{Verbatim}
\smallskip
\textbf{Possible sections:}
\begin{enumerate}
\setlength{\itemsep}{1pt}
\setlength{\parsep}{0pt}
\setlength{\parskip}{0pt}
\item Introduction: why Mizar needs evolution now.
\item Mizar as baseline: readability, MML, Formalized Mathematics.
\item Design principles and the three pillars.
\item Language evolution: dependencies, structures, registrations, templates.
\item Verifier architecture, certificates, and the small kernel.
\item Verified computation and the MVM.
\item AI-safe proof development.
\end{enumerate}
\note{
\smallskip
\textbf{Presenter script:}
\smallskip
\textbf{Read aloud:}
\begin{itemize}
\setlength{\itemsep}{1pt}
\setlength{\parsep}{0pt}
\setlength{\parskip}{0pt}
\item Introduction: why Mizar needs evolution now.
\item Mizar as baseline: readability, MML, Formalized Mathematics.
\item Design principles and the three pillars.
\end{itemize}
After the example, say which design obligation the syntax makes visible.

}
\end{frame}

\begin{frame}[fragile,t]{Backup D. Paper Outline Seed (2/2)}
\smallskip
\textbf{Possible sections:}
\begin{enumerate}
\setlength{\itemsep}{1pt}
\setlength{\parsep}{0pt}
\setlength{\parskip}{0pt}
\item Package-based library and publication workflow.
\item Migration plan and evaluation metrics.
\item Related work and collaboration agenda.
\end{enumerate}
\note{
\smallskip
\textbf{Presenter script:}
\smallskip
\textbf{Read aloud:}
\begin{itemize}
\setlength{\itemsep}{1pt}
\setlength{\parsep}{0pt}
\setlength{\parskip}{0pt}
\item Package-based library and publication workflow.
\item Migration plan and evaluation metrics.
\item Related work and collaboration agenda.
\end{itemize}
}
\end{frame}

\section{Backup E. Reviewer Checklist}
\begin{frame}[fragile,t]{Backup E. Reviewer Checklist}
\smallskip
\textbf{Use this checklist before converting to Beamer:}
\begin{itemize}
\setlength{\itemsep}{1pt}
\setlength{\parsep}{0pt}
\setlength{\parskip}{0pt}
\item Does every story open with a real cost, not a feature announcement?
\item Does every criticism acknowledge why the current practice was useful?
\item Does every code example carry its status label (exact MML excerpt, specification example, or sketch)?
\item Is every exact excerpt attributed with article and line numbers?
\item Does every story end with questions the audience can actually answer?
\item Are Red AI edits clearly forbidden?
\item Are migration claims measurable?
\end{itemize}
\smallskip
\textbf{Use this checklist before converting to Beamer:}
\begin{itemize}
\setlength{\itemsep}{1pt}
\setlength{\parsep}{0pt}
\setlength{\parskip}{0pt}
\item Is EBNF absent from all frames?
\end{itemize}
\note{
\smallskip
\textbf{Presenter script:}
\smallskip
\textbf{Read aloud:}
\begin{itemize}
\setlength{\itemsep}{1pt}
\setlength{\parsep}{0pt}
\setlength{\parskip}{0pt}
\item Does every story open with a real cost, not a feature announcement?
\item Does every criticism acknowledge why the current practice was useful?
\item Does every code example carry its status label (exact MML excerpt, specification example, or sketch)?
\item Is every exact excerpt attributed with article and line numbers?
\item Does every story end with questions the audience can actually answer?
\end{itemize}
}
\end{frame}

\end{document}
